%-----------------------------------------------------------------------------
% Rozdział V
%-----------------------------------------------------------------------------

\cleardoublepage
\chapter{Gomulscy z~Desna}

% Przednia okładka podrozdziału
\includepdf{Desno_mapa_fin.png}

\lettrine[lines=5, lhang=0.1]{D}{esno} jako obszar, jeszcze nie samodzielna
miejscowość, prawdopodobnie po~raz pierwszy zostało wymienione w~oficjalnym
dokumencie państwowym w~1772~roku. Dokumentem tym był \enquote{Inwentarz dóbr
dzierżawy Dembe} \cite{debe1}~sporządzony przez rotmistrza Andrzeja Olędzkiego
na polecenie Komisji Skarbowej Koronnej. Starostwo Dębe w~tamtym czasie było
częścią ziem królewskich, w~ziemi czerskiej województwa mazowieckiego,
które dzierżawił do 1772~roku podkanclerzy Wielkiego Księstwa Litewskiego
Antoni Tadeusz Przezdziecki herbu Roch III, po czym po jego śmierci zostało
ono przekazane w~dzierżawę jego synowi Dominikowi, późniejszemu Staroście
Mińskiemu.

\begin{figure}[!ht]
    \vspace*{0.5cm}
    \centering \includegraphics[width=1.0\linewidth]{
        Inwentarz_dobr_Debe_1.jpg}
    \captionsetup{format=hang}
    \caption{Fragment inwentarza dóbr dzierżawy Dębe z~1772~roku -- na
    czerwono podkreślono zdanie odnoszące się do boru \enquote{Deśnie}
    \cite{debe1}.}
    \label{fig:desna1_1772}
\end{figure}

Nazwa \enquote{Deśnie} została wymieniona już w~pierwszych zdaniach inwentarza
dóbr dzierżawy Dembe, opisujących granicę dóbr dzierżawy (patrz
ryc.\ref{fig:desna1_1772}):

\begin{centeredquote}
\textit{...około Pola Cisieńskiego do Dwuch Drog schodzących do Błota
nazwanego cięciowo od tej cięciwy do Bora głownego nazywaiącego się
\textbf{Deśnie} tąż Dróżką od Strugi Długa nazwana idący...}
\end{centeredquote}

Z~powyższego opisu wynika, iż bór znajdujący się nad rzeką Długą niedaleko pól
wsi Cisie i~bagien Cięciwa był już w~1772~roku znany pod nazwą
\enquote{Deśnie}\footnote{Aktualnie wieś Desno znajduje się w~odległości około
pół kilometra na północ od rzeki Długiej (przepływającej przez wieś Cisie)
oraz około dwóch kilometrów na zachód od wsi Cięciwa.}.

\begin{figure}[!ht]
    \vspace*{0.5cm}
    \centering \includegraphics[width=0.96\linewidth]{
        Inwentarz_dobr_Debe_2.jpg}
    \captionsetup{format=hang}
    \caption{Fragment inwentarza dóbr dzierżawy Dębe z~1772~roku -- na
    czerwono podkreślono zdanie odnoszące się do boru \emph{Desna}
    \cite{debe1}.}
    \label{fig:desna2_1772}
\end{figure}

Jednak to nie jedyny fragment inwentarza z~1772~roku, w~którym występuje
nazwa zbliżona do nazwy Desno. W~kolejnym fragmencie tego dokumentu
wymieniającym najważniejsze wsie, folwarki, pola, łąki oraz lasy składające
się na starostwo Dębę, można odnaleźć następujący wpis:
\enquote{\textbf{Desna} do Rzeki Długa nazwaney}  
(patrz ryc.\ref{fig:desna2_1772})\footnote{Wśród nazw wymienionych w~powyższym
fragmencie inwentarza, poza Desnem, można odnaleźć jeszcze wiele takich, które
do dziś funkcjonują jako nazwy okolicznych wsi, jak na przykład: Kobierne,
Rysie, Choszczówka, Cięciwa czy Cyganka.}. Powtórne wymienienie nazwy
\emph{Desna} w~tym dokumencie wskazuje, iż bór \emph{Deśnie} / \emph{Desna}
był w~1772~roku istotnym elementem składowym dóbr dzierżawy Dębe, a~sama jego
nazwa była już w~ówczesnym czasie ustalona i~powszechnie używana.

\begin{figure}[!ht]
    \vspace*{0.5cm}
    \centering \includegraphics[width=1.0\linewidth]{
        Inwentarz_dobr_Debe_3.jpg}
    \captionsetup{format=hang}
    \caption{Fragment dokumentu \enquote{Lustracya dzierżawy Dembego}
    z~1789~roku -- na czerwono podkreślono zdanie odnoszące się do lasu
    Desnego \cite{debe2}.}
    \label{fig:desna3_1789}
\end{figure}

W~kolejnym inwentarzu dóbr dzierżawy Dębe z~1782~roku, który został
sporządzony po~śmierci Dominika Przezdzieckiego, nie sposób jest odnaleźć
nazwy nawiązującej do omawianej w~niniejszym rozdziale miejscowości, ale
inwentarz ten jest zdecydowanie krótszy i~mniej szczegółowy od tego
z~1772~roku. Kolejny raz nazwa \emph{Desna} pojawia się w~inwentarzu dóbr
dzierżawy Dębe z~1789~roku, w~postaci następującego wpisu:

\begin{centeredquote}
\textit{...jako też Las \textbf{Desny} na Budowlę, Gonty, Tartaczną Fabrykę
użyteczny...}
\end{centeredquote}

Inwetarz ten w~tym fragmencie wskazuje, iż drzewo pozyskiwane z~lasu
\emph{Desnego} było wykorzystywane do produkcji budynków, gontów oraz do
obsługi tartaku (patrz ryc.\ref{fig:desna3_1789}). Mamy więc kolejne
potwierdzenie używania nazw \emph{Deśnie} / \emph{Desna} / \emph{Desny} już
w~końcówce XVIII~wieku w~odniesieniu do boru znajdującego się niedaleko rzeki
Długiej, będącego częścią starostwa Dębe, należącego do dóbr królewskich
Pierwszej Rzeczypospolitej. Skąd więc wzięła się nazwa tego boru?
Najprawdopodobniej pochodzi ona od słowa \enquote{desny}, które w~języku
staropolskim oznaczało \enquote{prawy} lub \enquote{po prawej stronie}. Bór
ten znajdował się bowiem po prawej stronie rzeki Długiej, patrząc w~kierunku
jej biegu i~był on częścią Puszczy Dębskiej, która rozciągała się na terenie
dzisiejszych gmin: Dębe Wielkie oraz Halinów.

\begin{figure}[!ht]
    \vspace*{0.5cm}
    \centering \includegraphics[width=0.83\linewidth]{
        Desna_ksiega_hipoteczna1.png}
    \captionsetup{format=hang}
    \caption{Pierwsze zdania księgi hipotecznej dóbr ziemskich \emph{Desna}
    z~1895~roku \cite{desna1}.}
    \label{fig:desna_hipoteka1}
\end{figure}

Nie zachowały się do dnia dzisiejszego żadne dalsze inwentarze dóbr dzierżawy
Dębe, w~których mogłaby pojawić się nazwa \emph{Desna}. Kolejnym, znanym
autorowi niniejszej książki, państwowym dokumentem, w~którym użyta została
nazwa \emph{Desna}, jest księga hipoteczna z~1895~roku o~nazwie \enquote{Dobra
ziemskie DESNA i~KĄTY MROWISKIE t.~1, powiat nowomiński, wł. Izrael
Rotgolc} \cite{desna1} należąca do zbioru \enquote{Hipoteka w Mińsku
Mazowieckim} Archiwum Państwowego w~Siedlcach. Księga ta spisana została
częściowo w~języku rosyjskim\footnote{W~dalszej części niniejszej książki, gdy
będę przywoływać daty w~kontekście wpisów w~języku rosyjskim, dla lepszej
czytelności będę stosować kalendarz gregoriański, pomimo iż wpisy te
oryginalnie datowane były zgodnie z~kalendarzem juliańskim, który obowiązywał
ówcześnie na terenie zaboru rosyjskiego.} (w~początkowej swojej części)
i~zawiera szczegółowy opis historii powstania miejscowości \emph{Desna}.
Pierwsze zdania księgi hipotecznej wsi \emph{Desna} wskazują, iż miejscowość
ta w~momencie powstania, czyli 11~lipca 1895~roku\footnote{Za datę powstania
wsi \emph{Desna} należy uznać datę pierwszego wpisu do księgi wieczystej, który
datowany jest na 29~czerwca 1895~roku (kalendarz juliański) / 11~lipca
1895~roku (kalendarz gregoriański) -- patrz ryc. \ref{fig:desna_hipoteka2}.
W~pierwszych zdaniach księgi hipotecznej pada wprawdzie data 13~lipca
1895~roku, czyli 25~lipca 1895~roku zgodnie z~kalendarzem gregoriańskim, ale
jest to prawdopodobnie data urzędowego zarejestrowania księgi hipotecznej po
jej utworzeniu.}, znajdowała się na terenie gminy Wiązowna, powiatu
nowomińskiego, w~obwodzie warszawskim i~liczyła sobie 144~dziesięciny oraz
1190~sążni kwadratowych powierzchni, czyli około 157~hektarów\footnote{Obecnie
wieś Desno również zajmuje około 157~hektarów powierzchni, co wskazywałoby, że
granice tej miejscowości nie uległy zbyt dużym zmianom w~ciągu 130~lat jej
historii.} (patrz ryc. \ref{fig:desna_hipoteka1}). Dalszy opis nowo powstałej
wsi \emph{Desna} w~księdze wieczystej wskazuje, iż od południa graniczyła ona
z~ziemiami folwarku prywatnego Cisie, od zachodu z~ziemiami chłopów wsi
Krzewina, od północnego zachodu z~ziemiami chłopów wsi Chobot, od północnego
wschodu z~ziemiami chłopów wsi Kąty Mrowiskie, a~od wschodu z~ziemiami
folwarku Olesinek.

\begin{figure}[!ht]
    \vspace*{0.5cm}
    \centering \includegraphics[width=0.86\linewidth]{
        Desna_ksiega_hipoteczna2.JPG}
    \captionsetup{format=hang}
    \caption{Data pierwszego wpisu do księgi wieczystej wsi \emph{Desna} oraz
    podpis ówczesnego właściciela tej wsi \cite{desna1}.}
    \label{fig:desna_hipoteka2}
\end{figure}

Z~dalszej treści księgi hipotecznej wynika, iż wieś \emph{Desna} została
wydzielona z~dóbr ziemskich sąsiedniej wsi Cisie\footnote{Co w~pełni pokrywa
się ze słowami babci autora niniejszej książki, Jadwigi Leokadii
Gomulskiej\indeksosoby[Poziemska]{Jadwiga Leokadia}{Gomulska}{1935}{Desno}
z~domu Poziemskiej, która twierdziła, iż jej matka, Magdalena
Poziemska\indeksosoby[Zawadzka][Gomulska]{Magdalena}{Poziemska}{1895}{Cisie}
z~domu Zawadzka, urodziła się w~\emph{Koloniach Cisieńskich} -- tak
prawdopodobnie zwyczajowo nazwano tę~część wsi Cisie przed powstaniem wsi
\emph{Desna}. Co~ważne, prababka autora niniejszej książki urodziła się
na~terenie obecnego Desna 25~czerwca 1895~roku, czyli na kilkanaście dni przed
oficjalnym powstaniem wsi \emph{Desna}.}, przez ówczesnego częściowego
właściciela tych dóbr Izraela Rotholca, mieszkańca Warszawy. Niedługo\
po~wydzieleniu wsi \emph{Desna} z~majątku Cisie -- 23~lipca 1895~roku, Izrael
Rotholc sprzedał ją w~całości założycielom spółki \emph{Desna} i~tym samym
stali się oni jedynymi właścicielami tych terenów. Był to ostatni etap
transakcji sprzedaży tych gruntów, która rozpoczęła się 5~kwietnia 1894~roku,
kiedy to~17~gospodarzy chłopskich założyło spółkę (towarzystwo) o~nazwie
\emph{Desna}. Już od~13~marca 1894~roku przyszli członkowie spółki mieli
możliwość osiedlania się na terenach przyszłej wsi \emph{Desna}, dzieląc się
od tego czasu po połowie, z~aktualnym właścicielem tych terenów Izraelem
Rotholcem, wszelkimi zobowiązaniami podatkowymi do momentu finalizacji
transakcji sprzedaży. Wraz z~założeniem spółki \emph{Desna} zawarta została
umowa, w~której jej członkowie zobowiązali się do~wspólnego nabycia
150~hektarów\footnote{Jak wcześniej wspominano, wieś \emph{Desna} liczyła
sobie około 157~hektarów, pozostałe 7~hektarów, jak wskazuje księga wieczysta,
stanowiły drogi i~granice.} ziemi rolnej wsi \emph{Desna} oraz zapłacili
Izraelowi Rotholcowi zadatek na~rzecz tej transakcji w~wysokości 1840~rubli
i~89~kopiejek.

\begin{figure}[!ht]
    \vspace*{0.5cm}
    \centering \includegraphics[width=0.84\linewidth]{
        Desna_zalozyciele.JPG}
    \captionsetup{format=hang}
    \caption{Tabela przedstawiająca 17 gospodarzy chłopskich, którzy byli
    członkami spółki \emph{Desna} w~momencie zakupu wsi \emph{Desna} od
    Izraela Rotholca 23~lipca~1895~roku \cite{desna1}.}
    \label{fig:desna_zalozyciele}
\end{figure}

Łączna cena sprzedaży \enquote{dóbr majątku \emph{Desna}} została ustalona
na~14405 rubli\footnote{Przeliczając tę kwotę na dzisiejsze pieniądze byłoby
to~około trzy miliony złotych.}. Udziałowcy spółki \emph{Desna}, 23~lipca
1895~roku, musieli więc zapłacić Izraelowi Rotholcowi pozostałą kwotę 12564
rubli i~11 kopiejek. Udało im się zebrać, do tego czasu, jedynie 4564~ruble
i~11~kopiejek, więc na~pozostałą kwotę 8000~rubli zaciągnęli wspólną pożyczkę
hipoteczną na~okres 26~lat i~pół roku w~warszawskim oddziale Banku Ziemskiego
Włościańskiego dla guberni warszawskiej\footnote{Od 1890~roku na terenie
zaboru rosyjskiego zaczęły funkcjonować pożyczki chłopskie na~preferencyjnych
warunkach, aby wzmocnić chłopską własność w~kontrze do~własności
szlacheckiej.}. Kwota pożyczki została rozdzielona na~udziałowców spółki
\emph{Desna} proporcjonalnie do~wielkości ziemi, jaką każdy z~nich nabywał
oraz skorygowana przez wielkość wpłaconych przez nich zadatków. I~tak
największa część pożyczki, bo~aż~1975~rubli przypadła Janowi Piotrkowiczowi,
a~najmniejsza część, bo~jedynie 375~rubli przypadła Janowi Jackowiczowi
(zięciowi Piotrkowicza). Pożyczka była spłacana w~półrocznych ratach
w~ustalonych przez bank kwotach z~solidarną wzajemną odpowiedzialnością
wszystkich udziałowców spółki \emph{Desna}. W~momencie finalizacji zakupu wsi
\emph{Desna} przez współzałożycieli spółki \emph{Desna}, jej członkami byli
następujący gospodarze chłopscy (patrz ryc. \ref{fig:desna_zalozyciele}):

\begin{enumerate}
    \item Jan Piotrkowicz\indeksosoby{Jan}{Piotrkowicz}{1847}{Kąty
    Goździejewskie} syn Tomasza (ur.~w~1847~r. w~Kątach
    Goździejewskich): nabył 19~dziesięcin i~518~sążni użytków rolnych oraz
    7~sążni nieużytków rolnych -- łącznie około 21~hektarów ziemi,
    \item Stanisław Wichrowski\indeksosoby{Stanisław}{Wichrowski}{1858}{Wólka
    Piecząca} syn Józefa (ur.~w~1858~r. w~Wólce Pieczącej pod
    Stanisławowem): nabył 15~dziesięcin i~835~sążni użytków rolnych oraz
    65~sążni nieużytków rolnych -- łącznie około 17~hektarów ziemi,
    \item Franciszek Rogala\indeksosoby{Franciszek}{Rogala}{1848}{Kąty
    Goździejewskie} syn Tomasza (ur.~w~1848~r. w~Kątach
    Goździejewskich): nabył 5~dziesięcin i~1767~sążni użytków rolnych oraz
    70~sążni nieużytków rolnych -- łącznie około 6~hektarów ziemi,
    \item Adam Ludwiniak\indeksosoby{Adam}{Ludwiniak}{1847}{Wola
    Cygowska} syn Antoniego (ur.~w~1847~r. w~Woli Cygowskiej pod
    Poświętnem): nabył 5~dziesięcin i~1808~sążni użytków rolnych oraz
    30~sążni nieużytków rolnych -- łącznie około 6~hektarów ziemi,
    \item Franciszek Bok\indeksosoby{Franciszek}{Bok}{1849}{Brzeziny} syn
    Wojciecha (ur.~w~1849~r. w~Brzezinach): nabył
    17~dziesięcin i~1850~sążni użytków rolnych oraz 1~dziesięcinę
    i~1075~sążni nieużytków rolnych -- łącznie około 21~hektarów ziemi,
    \item Jan Jackowicz\indeksosoby{Jan}{Jackowicz}{1869}{Poręby} syn
    Franciszka (ur.~w~1869~r. w~Porębach): nabył
    3~dziesięciny i~1995~sążni użytków rolnych oraz 30~sążni nieużytków
    rolnych -- łącznie około 4~hektary ziemi,
    \item Konstanty Kaim\indeksosoby{Konstanty}{Kaim}{1860}{Stanisławów} syn
    Teodora (ur.~w~1860~r. w~Stanisławowie): nabył
    7~dziesięcin i~1413~sążni użytków rolnych oraz 237~sążni nieużytków
    rolnych -- łącznie około 8~hektarów ziemi,
    \item Franciszek Kutrowski\indeksosoby{Franciszek}{Kutrowski}{1871}{Dąbrówka
    Grzybowska} syn Jakuba (ur.~w~1871~r. w~Dąbrówce Grzybowskej
    pod Warszawą): nabył 5~dziesięcin i~1789~sążni użytków rolnych oraz
    48~sążni nieużytków rolnych -- łącznie około 6~hektarów ziemi,
    \item Bronisław Piotrowski\indeksosoby{Bronisław}{Piotrowski}{1868}{Wola
    Rafałowska} syn Piotra (ur.~w~1868~r. w~Woli Rafałowskiej
    pod Mrozami): nabył 5~dziesięcin i~1792~sążni użytków rolnych oraz
    46~sążni nieużytków rolnych -- łącznie około 6~hektarów ziemi,
    \item Franciszek Tkaczyk\indeksosoby{Franciszek}{Tkaczyk}{1862}{Olesin}
    syn Piotra (ur.~w~1862~r. w~Kątach Mrowiskich): nabył 7~dziesięcin
    i~1579~sążni użytków rolnych oraz 71~sążni nieużytków rolnych -- łącznie
    około 8~hektarów ziemi,
    \item Marcin Zawadzki\indeksosoby{Marcin}{Zawadzki}{1856}{Anielina} syn
    Józefa (ur.~w~1856~r. w~Anielinie pod Mińskiem Mazowieckim): nabył
    7~dziesięcin i~556~sążni użytków rolnych oraz
    1094~sążni nieużytków rolnych -- łącznie około 8~hektarów ziemi,
    \item Stanisław Łabędź\indeksosoby{Stanisław}{Łabędź}{1840}{Michałów}[Lxyzabedz]
    syn Mikołaja (ur.~w~1840~r. w~Michałowie pod Okuniewem): nabył
    6~dziesięcin i~2023~sążni użytków rolnych oraz 2027~sążni nieużytków
    rolnych -- łącznie około 8~hektarów ziemi,
    \item Stanisław Borucki\indeksosoby{Stanisław}{Borucki}{1858}{Łęka}[][Lxyzeka]
    syn Walentego (ur.~w~1858~r. w~Łęce pod Poświętnem): nabył 3~dziesięciny
    i~1591~sążni użytków rolnych oraz
    1259~sążni nieużytków rolnych -- łącznie około 5~hektarów ziemi,
    \item Marcin Gruba\indeksosoby{Marcin}{Gruba}{1858}{Wólka Barcząca} syn
    Pawła (ur.~w~1858~r. w~Wólce Barczącej pod Mińskiem Mazowieckim): nabył
    5~dziesięcin i~646~sążni użytków rolnych
    oraz 1289~sążni nieużytków rolnych -- łącznie około 6~hektarów ziemi,
    \item Michał Gomulski\indeksosoby{Michał}{Gomulski}{1859}{Anielina} syn
    Piotra (ur.~w~1859~r. w~Anielinie): nabył
    6~dziesięcin i~2158~sążni użytków rolnych oraz 1892~sążni nieużytków --
    łącznie około 8~hektarów ziemi,
    \item Stanisław Gomulski\indeksosoby{Stanisław}{Gomulski}{1863}{Anielina}
    syn Piotra (ur.~w~1863~r. w~Anielinie): nabył
    6~dziesięcin i~1304~sążni użytków rolnych oraz 1~dziesięcinę i~346~sążni
    nieużytków rolnych -- łącznie około 8~hektarów ziemi,
    \item Antoni Zgutka\indeksosoby{Antoni}{Zgutka}{1858}{Arynów} syn
    Franciszka (ur.~w~1858~r. w~Arynowie pod Mińskiem Mazowieckim): nabył
    3~dziesięciny i~1353~sążni użytków rolnych oraz 672~sążni nieużytków
    rolnych -- łącznie około 4~hektary ziemi.
\end{enumerate}

Większość założycieli spółki \emph{Desna} pochodziła z~miejscowości oddalonych
od wsi \emph{Desna} o~kilkanaście lub kilkadziesiąt kilometrów -- zaledwie
sześciu z~nich można nazwać okolicznymi gospodarzami (z~Kątów Goździejewskich,
Kątów Mrowiskich, Poręb, Brzezin czy Michałowa). Jak wynika z~analizy drzew
genealogicznych założycieli spółki \emph{Desna}, wielu z~nich było ze~sobą
powiązanych więzami rodzinnymi zanim stali się sąsiadami, stąd prawdopodobnie
łatwiej było im~zaciągnąć wspólne długoletnie zobowiązanie i~zamieszkać obok
siebie na~terenie nowo powstałej wsi:

\begin{itemize}
    \item Żony Jana Piotrkowicza\indeksosoby{Jan}{Piotrkowicz}{1847}{Kąty
    Goździejewskie} oraz Franciszka
    Boka\indeksosoby{Franciszek}{Bok}{1849}{Brzeziny} były siostrami -- żoną
    Jana była Krystyna\indeksosoby[Link]{Krystyna}{Piotrkowicz}{1845}{Olesin}
    z~domu Link (ur.~w~1845~r. w~Kątach Mrowiskich), a~żoną
    Franciszka była Teofila\indeksosoby[Link]{Teofila}{Piotrkowicz}{1848}{Olesin}
    z~domu Link (ur.~w~1848~r. w~Kątach Mrowiskich),
    \item Żoną Stanisława Wichrowskiego\indeksosoby{Stanisław}{Wichrowski}{1858}{Wólka
    Piecząca} była Józefa\indeksosoby[Piotrkowicz]{Józefa}{Wichrowska}{1859}{Kąty
    Goździejewskie} z~domu Piotrkowicz (ur.~w~1859~r. w~Kątach Goździejewskich)
    siostra Jan Piotrkowicza,
    \item Żoną Jana Jackowicza\indeksosoby{Jan}{Jackowicz}{1869}{Poręby} była
    Franciszka\indeksosoby[Piotrkowicz]{Franciszka}{Jackowicz}{1879}{Olesin}
    z~domu Piotrkowicz (ur.~w~1879~r. w~Olesinie) córka
    Jana\indeksosoby{Jan}{Piotrkowicz}{1847}{Kąty Goździejewskie}
    i~Krystyny\indeksosoby[Link]{Krystyna}{Piotrkowicz}{1845}{Olesin}
    Piotrkowiczów,
    \item Żoną Franciszka Kutrowskiego\indeksosoby{Franciszek}{Kutrowski}{1871}{Dąbrówka
    Grzybowska} była Marianna Paulina\indeksosoby[Piotrowska]{Marianna
    Paulina}{Kutrowska}{1875}{Wola Rafałowska} z~domu Piotrowska
    (ur.~w~1875~r. w~Woli Rafałowskiej pod Mrozami) siostra
    Bronisława Piotrowskiego\indeksosoby{Bronisław}{Piotrowski}{1868}{Wola
    Rafałowska},
    \item Żoną Stanisława Boruckiego\indeksosoby{Stanisław}{Borucki}{1858}{Łęka}[][Lxyzeka]
    była Marianna\indeksosoby[Ludwiniak]{Marianna}{Borucka}{1845}{Wola
    Cygowska} z~domu Ludwiniak (ur.~w~1845~r. w~Woli Cygowskiej pod Poświętnem)
    siostra Adama Ludwiniaka\indeksosoby{Adam}{Ludwiniak}{1847}{Wola
    Cygowska},
    \item Żoną Marcina Gruby\indeksosoby{Marcin}{Gruba}{1858}{Wólka Barcząca}
    była Marianna\indeksosoby[Zawadzka]{Marianna}{Gruba}{1867}{Anielina}
    z~domu Zawadzka (ur.~w~1867~r. w~Anielinie) siostra Marcina
    Zawadzkiego\indeksosoby{Marcin}{Zawadzki}{1856}{Anielina}.
    \item Gomulscy i~Zawadzcy nie tylko pochodzili z~jednej wsi pod Mińskiem
    Mazowieckim, ale również byli rodzicami chrzestnymi swoich dzieci.
\end{itemize}

\begin{figure}[!ht]
    \vspace*{0.5cm}
    \centering \includegraphics[width=0.82\linewidth]{
        Desna_zalozyciele2.jpg}
    \captionsetup{format=hang}
    \caption{Tabela przedstawiająca właścicieli ziemskich we wsi \emph{Desna} 
    w~1914~roku \cite{desna1}.}
    \label{fig:desna_zalozyciele2}
\end{figure}

Nie wszyscy założyciele spółki \emph{Desna} zadomowili się na stałe na terenie
nowo powstałej wsi\footnote{Niektórzy nigdy tam nie mieszkali, mieli tam tylko
pola uprawne -- jak np.~największy udziałowiec spółki \emph{Desna} Jan
Piotrkowicz\indeksosoby{Jan}{Piotrkowicz}{1847}{Kąty Goździejewskie}, który
mieszkał w~sąsiednim Olesinie Górnym.} -- w~ciągu pierwszych 20~lat istnienia
miejscowości \emph{Desna} część z~nich sprzedała swoje udziały lub przeniosła
je~na~swoje dzieci (patrz ryc. \ref{fig:desna_zalozyciele2}):

\begin{itemize}
    \item Pierwszymi nowymi mieszkańcami wsi \emph{Desna} była rodzina
    Stanisława Gomulskiego\indeksosoby{Stanisław}{Gomulski}{1861}{Wólka Mińska}
    (ur.~1861~r.), syna Jana i~Agnieszki\footnote{Co
    ciekawe matka Stanisława, Agnieszka z~domu Piotrkowicz, była dalszą
    kuzynką Jana Piotrkowicza.}, a~bratanka Michała oraz Stanisława (już
    mieszkających na Deśnie). Około 1897~roku Stanisław przeniósł się wraz
    z~rodziną do nowo powstałej wsi (świadczy o~tym akt chrztu jego syna
    Franciszka), a~6~października 1900~roku odkupił finalnie udziały w~spółce
    \emph{Desna} za kwotę 600~rubli od Konstatntego
    Kaima\indeksosoby{Konstanty}{Kaim}{1860}{Stanisławów} przejmując w~ten
    sposób własność 8~hektarów ziemi rolnej na~terenie wsi \emph{Desna} oraz
    dług na kwotę około 500~rubli.
    \item Na początku 1901~roku za kwotę 586~rubli i~63~kopiejek Michał
    Gomulski\indeksosoby{Michał}{Gomulski}{1859}{Anielina} odkupił udziały
    w~spółce \emph{Desna} od Marcina Gruby\indeksosoby{Marcin}{Gruba}{1858}{Wólka
    Barcząca}, przejmując w~ten sposób własność kolejnych 6~hektarów ziemi.
    Tym samym Michał Gomulski stał się czwartym co~do~wielkości właścicielem
    ziemskim na~terenie wsi \emph{Desna} -- posiadał już bowiem 14~hektarów
    ziemi, pod tym względem ustępując jedynie Janowi
    Piotrkowiczowi\indeksosoby{Jan}{Piotrkowicz}{1847}{Kąty Goździejewskie}
    (21~hektarów), Franciszkowi
    Bokowi\indeksosoby{Franciszek}{Bok}{1849}{Brzeziny} (21~hektarów) oraz
    Stanisławowi Wichrowskiemu\indeksosoby{Stanisław}{Wichrowski}{1858}{Wólka
    Piecząca} (17~hektarów).
    \item 13 marca 1903~roku Antoni
    Poziemski\indeksosoby{Antoni}{Poziemski}{1864}{Cisie} (ur.~w~1864~roku
    w~Cisiu) odkupił udziały w~spółce \emph{Desna} od~Bronisława
    Piotrowskiego\indeksosoby{Bronisław}{Piotrowski}{1868}{Wola Rafałowska}
    za~kwotę 400~rubli przejmując w~ten sposób własność około 6~hektarów ziemi
    oraz przypisany do~niej dług. Żona Antoniego Poziemskiego,
    Józefa\indeksosoby[Tomaszewska]{Józefa}{Poziemska}{1864}{Chobot} z~domu
    Tomaszewska (ur.~w~1864 roku w~Chobocie), była starszą siostrą
    Franciszki\indeksosoby[Tomaszewska]{Franciszka}{Piotrowska}{1868}{Chobot}
    (ur.~w~1868~roku w~Chobocie) żony Bronisława Piotrowskiego. 
    \item Około 1904~roku Jan
    Piotrkowicz\indeksosoby{Jan}{Piotrkowicz}{1847}{Kąty Goździejewskie}
    przekazał większość swoichnudziałów w~spółce \emph{Desna} swojemu synowi
    Andrzejowi\indeksosoby{Andrzej}{Piotrkowicz}{1885}{Olesin} (ur.~w~1885~roku
    w~Olesinie) i~jego żonie
    Magdalenie\indeksosoby[Malesa]{Magdalena}{Piotrkowicz}{1885}{Wielgolas
    Brzeziński} z~domu Malesa (ur.~w~1885~roku w~Wielgolesie Brzezińskim).
    \item Andrzej Karłowski\indeksosoby{Andrzej}{Karłowski}{1877}{Cisie}
    (ur.~w~1877~roku w~Cisiu) około 1910~roku odkupił udziały w~spółce
    \emph{Desna} od Antoniego Zgutki\indeksosoby{Antoni}{Zgutka}{1858}{Arynów}
    przejmując w~ten sposób własność około czterech hektarów ziemi,
    na~zachodnim krańcu wsi.
    \item W~1910~rok Szymon
    Rogala\indeksosoby{Szymon}{Rogala}{1878}{Kąty Goździejewskie}
    (ur.~w~1878~roku w~Kątach Goździejewskich), syn
    Franciszka\indeksosoby{Franciszek}{Rogala}{1848}{Kąty Goździejewskie},
    wziął ślub z~Zofią
    Ludwiniak\indeksosoby[Ludwiniak]{Zofia}{Rogala}{1894}{Chobot}
    (ur.~w~1894~roku w~Chobocie), córką Adama, a~ich ojcowie przepisali
    na~nich swoje udziały w~spółce \emph{Desna} -- w~ten sposób
    Szymon i~Zofia stali się właścicielami około 12~hektarów ziemi rolnej
    w~północnej części wsi,
    \item Około 1912~roku Stanisław
    Gomulski\indeksosoby{Stanisław}{Gomulski}{1863}{Anielina} (syn Piotra)
    przekazał większość swoich udziałów w~spółce \emph{Desna} swojemu synowi
    Janowi\indeksosoby{Jan}{Gomulski}{1882}{Jarząbek} (ur.~w~1882~r. we wsi
    Jarząbek pod Mińskiem Mazowieckim).
    \item Franciszek Bok\indeksosoby{Franciszek}{Bok}{1849}{Brzeziny} około
    1914~roku przekazał część swoich udziałów w~spółce \emph{Desna} swojemu
    synowi Tomaszowi\indeksosoby{Tomasz}{Bok}{1873}{Aleksandrówka}
    (ur.~w~1873~roku w~Aleksandrówce).
    \item Michał Gomulski\indeksosoby{Michał}{Gomulski}{1859}{Anielina} około
    1914~roku przekazał ziemię odkupioną od Marcina
    Gruby\indeksosoby{Marcin}{Gruba}{1858}{Wólka Barcząca} oraz związane z~nią
    udziały w~spółce \emph{Desna} swojemu synowi
    Janowi\indeksosoby{Jan}{Gomulski}{1888}{Mikanów} (ur.~w~1888~r. w~Mikanowie).
\end{itemize}

\begin{figure}[!ht]
    \vspace*{0.5cm}
    \centering \includegraphics[width=1.05\linewidth]{
        Desna_1915.jpg}
    \captionsetup{format=hang}
    \caption{Mapa przedstawiająca orientacyjny podział wsi \emph{Desna}
    na poszczególne rodziny w~1915 roku.}
    \label{fig:desna_1915}
\end{figure}

Mapa na ryc. \ref{fig:desna_1915} przedstawia orientacyjny podział wsi
\emph{Desna} na poszczególne rodziny w~1915~roku, czyli po około 20~latach
istnienia tej miejscowości. Oznaczenie:
\begin{itemize}
\item \emph{Gomulscy 1} odnosi się do rodziny
Michała\indeksosoby{Michał}{Gomulski}{1859}{Anielina}
i~Marianny\indeksosoby[Wróbel]{Marianna}{Gomulska}{1860}{Siennica} Gomulskich,
\item \emph{Gomulscy 2} do rodziny
Stanisława\indeksosoby{Stanisław}{Gomulski}{1863}{Anielina}
i~Elżbiety\indeksosoby[Brodzińska]{Elżbieta}{Gomulska}{1857}{Grębiszew}
Gomulskich
\item \emph{Gomulscy 3} do rodziny
Stanisława\indeksosoby{Stanisław}{Gomulski}{1861}{Wólka Mińska}
i~Julianny\indeksosoby[Gągol]{Julianna}{Gomulska}{1862}{Dłużka}
Gomulskich\footnote{Numeracja rodzin Gomulskich jest zgodna z~kolejnością ich
osiedlania się na terenie wsi \emph{Desna}}.
\end{itemize}
Jak wynika z~przedstawionej mapy, w~1915~roku rodziny Gomulskich posiadały
łącznie około 30~hektarów ziemi rolnej we wsi \emph{Desna}, co~stanowiło
niemal 20\% powierzchni całej wsi.

\begin{figure}[!ht]
    \vspace*{0.5cm}
    \centering \includegraphics[width=0.9\linewidth]{
        1895_Katarzyna_Gomulska_Zawadzka_akt_chrztu_parafia_Długa_Kościelna_wpis_108.jpg}
    \captionsetup{format=hang}
    \caption{Akt chrztu Katarzyny Gomulskiej -- par. Długa Kościelna 1895~rok
    (108/1895) \cite{par_dluga}.}
    \label{fig:kgomulska2_1895}
\end{figure}

Wieś \emph{Desna}, od razu po utworzeniu, została włączona do parafii świętej
Anny w~Długiej Kościelnej. Pierwszy wpis w~księgach metrykalnych tej parafii,
w~którym trzykrotnie pada nazwa \emph{Desna}, pochodzi z~listopada 1895~roku
-- jest to akt chrztu Katarzyny Gomulskiej córki Stanisława
i~Elżbiety Gomulskich (patrz ryc. \ref{fig:kgomulska2_1895}).

\begin{figure}[!ht]
    \vspace*{0.5cm}
    \centering \includegraphics[width=0.88\linewidth]{
        Desna_uchwala.jpg}
    \captionsetup{format=hang}
    \caption{Uchwała spółki \emph{Desna} z~1918~roku \cite{desna1}.}
    \label{fig:desna_uchwala}
\end{figure}

\begin{figure}[!ht]
    \vspace*{0.5cm}
    \centering \includegraphics[width=0.78\linewidth]{
        Desna_dzialka_13_splata.jpg}
    \captionsetup{format=hang}
    \caption{Zaświadczenie z~dnia 23~stycznia 1919~roku o~spłacie pożyczki
    zaciągniętej dla działki nr~13 w~miejscowości \emph{Desna} wydane przez
    Inspektorat Skarbowy w~powiecie warszawskim \cite{desna1}.}
    \label{fig:desna_dzialka_13}
\end{figure}

Każda zmiana właściciela ziemskiego na terenie wsi \emph{Desna} wymagała
odpowiedniej uchwały spółki \emph{Desna}, zgody banku oraz wpisu do księgi
hipotecznej. Przykładową uchwałę z~1918 roku w~sprawie wykreślenia z~listy
członków spółki \emph{Desna} Stanisława
Gomulskiego\indeksosoby{Stanisław}{Gomulski}{1863}{Anielina} (syna Piotra)
oraz wpisania na listę członków jego syna
Józefa\indeksosoby{Józef}{Gomulski}{1892}{Anielina} można zobaczyć na ryc. 
\ref{fig:desna_uchwala}.

\begin{figure}[!ht]
    \vspace*{0.5cm}
    \centering \includegraphics[width=0.88\linewidth]{Desna_koniec.jpg}
    \captionsetup{format=hang}
    \caption{Zaświadczenie o~całkowitej spłacie pożyczki zaciągniętej przez
    członków spółki \emph{Desna} wydane przez Państwowy Bank Rolny
    w~Warszawie \cite{desna1}.}
    \label{fig:desna_koniec}
\end{figure}

Kolejne wpisy w~księdze hipotecznej wsi \emph{Desna} nie przynoszą już
żadnych nowych informacji na temat historii tej miejscowości -- jedynie
potwierdzają one kolejne zmiany właścicieli gruntów na jej terenie. Wśród
tych wpisów można znaleźć nowe rodziny, które stały się właścicielami
ziemskimi na~terenie wsi \emph{Desna} po~1915~roku, a~ich potomkowie żyją tam
do~dnia dzisiejszego. Należą do~nich między innymi:

\begin{itemize}
    \item Rodzina Parkotów: Julia Bok\indeksosoby[Bok]{Julia}{Parkot}{1881}{Dębe
    Wielkie} (ur.~w~1881~r. w~Dębem Wielkim) w~1905~roku wyszła za~mąż
    za~Pawła Parkota\indeksosoby{Paweł}{Parkot}{1883}{Glinianka} (ur.~w~1883~r.
    w~Gliniance) a~około 1922~roku odziedziczyli oni około sześć hektarów ziemi
    rolnej po~Franciszku Boku\indeksosoby{Franciszek}{Bok}{1849}{Brzeziny}.
    \item Rodzina Dworzników: Marianna
    Łabędź\indeksosoby[Łabędź]{Marinna}{Dworznik}{1880}{Brzeziny}[Lxyzabedz]
    (ur.~w~1880~r. w~Brzezinach) w~1901~roku wyszła za mąż za 
    Jana Dworznika\indeksosoby{Jan}{Dworznik}{1879}{Długa Szlachecka}
    (ur.~w~1879~r. w~Długiej Szlacheckiej) a~w~1926~roku odziedziczyli oni
    około cztery hektary ziemi rolnej po Stanisławie
    Łabędziu\indeksosoby{Stanisław}{Łabędź}{1840}{Michałów}[Lxyzabedz].
    \item Rodzina Lesińskich: Aniela
    Łabędź\indeksosoby[Łabędź]{Aniela}{Lesińska}{1875}{Brzeziny}[Lxyzabedz]
    (ur.~w~1875~r. w~Brzezinach) w~1910~roku wyszła za~mąż za~Jakuba
    Lesińskiego\indeksosoby{Jakub}{Lesiński}{1882}{Olesin} (ur.~w~1882~r.
    w~Olesinie) a~w~1926~roku odziedziczyli oni około cztery hektary ziemi
    rolnej po Stanisławie Łabędziu\indeksosoby{Stanisław}{Łabędź}{1840}{Michałów}[Lxyzabedz].
    \item Rodzina Krawczyńskich: Janina
    Kutrowska\indeksosoby[Kutrowska]{Janinna}{Krawczyńska}{1899}{Desno}
    (ur.~w~1899~r. na Deśnie) w~1918~roku wyszła za~mąż za~Stanisława
    Krawczyńskiego\indeksosoby{Stanisław}{Krawczyński}{1896}{Olesin}
    (ur.~w~1896~r. w~Olesinie) a~około 1930~roku odziedziczyli oni część
    ziemi rolnej po~Franciszku
    Kutrowskim\indeksosoby{Franciszek}{Kutrowski}{1871}{Dąbrówka Grzybowska}.
    \item Rodzina Zawadów: Wiktoria
    Gomulska\indeksosoby[Gomulska]{Wiktoria}{Zawada}{1916}{Desno}
    (ur.~w~1916~r. na Deśnie) w~1937~roku wyszła za~mąż za~Czesława Mariana
    Zawadę\indeksosoby{Czesław Marian}{Zawada}{1912}{Duchnów} (ur.~w~1912~r.
    w~Duchnowie) a~po kilku latach odziedziczyli oni część ziemi rolnej
    po~Janie Gomulskim\indeksosoby{Jan}{Gomulski}{1888}{Mikanów} (synu
    Stanisława i~Elżbiety Gomulskich).
    \item Rodzina Matuszewskich: Janina
    Gomulska\indeksosoby[Gomulska]{Janina}{Matuszewska}{1923}{Desno}
    (ur.~w~1923~r. na Deśnie) w~1941~roku wyszła za~mąż
    za~Czesława Matuszewskiego\indeksosoby{Czesław}{Matuszewski}{1917}{Cisówka}
    (ur.~w~1917~r. w~Cisówce) a~po kilku latach odziedziczyli oni część ziemi
    rolnej po Janie Gomulskim (synu Stanisława i~Elżbiety Gomulskich).
\end{itemize}

Spółka \emph{Desna} zakończyła swoją działalność w~lipcu 1922~roku, kiedy to
Państwowy Bank Rolny w~Warszawie (następca prawny Banku Ziemskiego
Włościańskiego dla guberni warszawskiej) potwierdził całkowitą spłatę pożyczki
zaciągniętej przez członków spółki w~1895~roku i~nakazał wykreślenie wpisu
dotyczącego hipoteki z~księgi wieczystej wsi \emph{Desna} (patrz ryc.
\ref{fig:desna_koniec}). Po rozwiązaniu spółki \emph{Desna} nazwa wsi
\emph{Desna} zaczęła stopniowo zanikać, zarówno w~oficjalnych dokumentach
państwowych, jak~i w~języku mówionym\footnote{Do tego stopnia, że autor
niniejszej książki, wychowany w~latach 90.~XX~wieku na Deśnie, przez potomków
kilku założycieli spółki \emph{Desna}, pierwszy raz spotkał się z~nazwą
\emph{Desna} analizując stare mapy województwa mazowieckiego}. Nazwę tę
zaczęła wypierać nazwa Desno, która była bardziej zgodna z~polską fonetyką.
W~dokumentach urzędowych czy mapach (patrz ryc. \ref{fig:desna_1925_1937})
nazwa \emph{Desna} pojawiała się jeszcze sporadycznie aż do lat 60.~XX wieku.
Pierwszy oficjalny dokument państwowy, który udało się odnaleźć autorowi
niniejszej książki, a~w~którym występuje nazwa Desno, pochodzi z~3~czerwca
1921~roku i~jest to projekt administracyjnego podziału gmin powiatu
mińskiego\footnote{Zgodnie z~najlepszą wiedzą autora niniejszej
książki, projekt ten nigdy nie wszedł w~życie i~wieś \emph{Desna} / Desno nie
była nigdy częścią gminy Dębę Wielkie} opublikowany w~Gazecie Urzędowej
Wydziału Powiatowego Starostwa Mińsko-Mazowieckiego dnia 20~lipca 1921~roku
(patrz ryc. \ref{fig:powiat_1921}). Ostatni oficjalny dokument państwowy,
który udało się odnaleźć autorowi niniejszej książki, a~w~którym występuje
nazwa \emph{Desna}, pochodzi natomiast z~31~maja 1960~roku -- jest to wykaz
urzędów stanu cywilnego w~powiecie Otwock opublikowany w~dzienniku urzędowym
Wojewódzkiej Rady Narodowej w~Warszawie (patrz ryc. \ref{fig:desna_1960}).

\begin{figure}[!ht]
    \vspace*{0.5cm}
    \centering \includegraphics[width=0.9\linewidth]{
        Powiat_1921.png}
    \captionsetup{format=hang}
    \caption{Fragment projektu administracyjnego podziału gmin powiatu
    Mińsko-Mazowieckiego z~dnia 3~czerwca 1921~roku \cite{gazeta_urzedowa}.}
    \label{fig:powiat_1921}
\end{figure}

\begin{figure}[!ht]
    \centering
    % Pierwsza mapa
    \includegraphics[width=0.95\linewidth]{Desna_mapa_1925.jpg}
    
    \vspace{0.1cm} % Odstęp między obrazami

    % Druga mapa
    \includegraphics[width=0.95\linewidth]{Desno_mapa_1937.jpg}
    
    % Wspólny podpis dla obu map
    \captionsetup{format=hang}
    \caption{Fragmenty map \enquote{Mapa Taktyczna Polski} z~1925~roku oraz
    1937~roku przedstawiające okolice wsi \emph{Desna} / Desno.}
    \label{fig:desna_1925_1937}
\end{figure}

\begin{figure}[!ht]
    \vspace*{0.5cm}
    \centering \includegraphics[width=1.0\linewidth]{
        Desna_1960.png}
    \captionsetup{format=hang}
    \caption{Fragment wykazu urzędów stanu cywilnego w~powiecie Otwock
    opublikowany w~Dzienniku Urzędowym Wojewódzkiej Rady Narodowej
    w~Warszawie dnia 31~maja 1960~roku \cite{dziennik_urzedowy}.}
    \label{fig:desna_1960}
\end{figure}

We wrześniu 1921~roku przeprowadzono na terenie Polski pierwszy powszechny
spis ludności po odzyskaniu niepodległości. Spis ten wykazał, iż na terenie
wsi \emph{Desna} mieszkało wówczas 116~osób\footnote{Dla porównania, równo
100~lat później, w~2021 roku, przeprowadzono spis powszechny, który wykazał,
iż~Desno miało w~tamtym momencie 229~mieszkańców.} (61~kobiet i~55~mężczyzn)
o~narodowości polskiej, zamieszkujących 18~budynków mieszkalnych (patrz ryc.
\ref{fig:desna_spis}).

\begin{figure}[!ht]
    \vspace*{0.5cm}
    \centering \includegraphics[width=1.0\linewidth]{
        Desna_spis_powszechny_1921.png}
    \captionsetup{format=hang}
    \caption{Fragment tabeli przedstawiającej statystki ludności
    zamieszkującej poszczególne wsie gminy Wiązowna w~1921~roku
    \cite{spis1921}.}
    \label{fig:desna_spis}
\end{figure}

Jak widać na ryc. \ref{fig:desna_spis} wieś \emph{Desna} została w~trakcie
spisu zakwalifikowana jako cześć gminy Wiązowna w~powiecie warszawskim, co
nie było do końca prawdą, gdyż przeniesienie gminy Wiązowna z~powiatu
Mińsko-Mazowieckiego do powiatu warszawskiego zostało prawnie uregulowane
przez władze II Rzeczpospolitej dopiero w~1924~roku (patrz ryc.
\ref{fig:wiazowna_warszawski}). Skąd więc taka nieścisłość w~spisie
powszechnym z~1921~roku? Z~powodu faktu, iż~gmina Wiązowna została wcielona do
powiatu warszawskiego już wcześniej -- w~1916~roku przez okupacyjne władze
niemieckie, ale decyzja ta nie została prawnie zatwierdzona przez władze II
Rzeczpospolitej aż do 1924~roku. Co oznacza, że w~1921~roku, gmina Wiązowna
\emph{de iure} należała do powiatu mińskiego, ale \emph{de facto} była już
zarządzana jako część powiatu warszawskiego.

\begin{figure}[!ht]
    \vspace*{0.5cm}
    \centering \includegraphics[width=0.6\linewidth]{
        Wiazowna_warszawski.png}
    \captionsetup{format=hang}
    \caption{Fragment rozporządzenia Rady Ministrów z~dnia 2~lipca 1924~roku
    o~uregulowaniu pod względem prawnym granic powiatu warszawskiego
    \cite{dziennik_ustaw3}.}
    \label{fig:wiazowna_warszawski}
\end{figure}

\begin{figure}[!ht]
    \vspace*{0.5cm}
    \centering \includegraphics[width=0.6\linewidth]{
        Desna_gmina_Okuniew.png}
    \captionsetup{format=hang}
    \caption{Fragment rozporządzenia Ministra Spraw Wewnętrznych generała
    Sławoja Składkowskiego z~dnia 31~lipca 1929~roku \cite{dziennik_ustaw}.}
    \label{fig:okuniew_1930}
\end{figure}

Pierwszego kwietnia 1930~roku, w~wyniku rozporządzenia Ministra Spraw
Wewnętrznych generała Sławoja Składkowskiego, z~dnia 31~lipca~1929~roku, wieś
\emph{Desna} po blisko 35~latach istnienia, wraz z~kilkoma okolicznymi
miejscowościami, została wydzielona z~gminy Wiązowna i~przeniesiona do gminy
Okuniew -- pozostając w~powiecie i~województwie warszawskim (patrz ryc.
\ref{fig:okuniew_1930})\footnote{Zauważmy, że w~rozporządzeniu nie ma mowy
o~wsiach Chobot i~Mrowiska, które znajdowały się zdecydowanie bliżej Okuniewa
niż Desno -- Chobot i~Mrowiska do 1952~roku były częścią gminy Dębe Wielkie,
zgodnie z~ich historyczną przynależnością do starostwa Dembe, a~potem razem
z~Desnem stały się częścią nowo powstałej gminy Halinów.}.

\begin{figure}[!ht]
    \vspace*{0.5cm}
    \centering \includegraphics[width=0.85\linewidth]{
        Spis_powszechny_1943.png}
    \captionsetup{format=hang}
    \caption{Fragment tabeli przedstawiającej statystki ludności
    zamieszkującej poszczególne miejscowości dystryktu warszawskiego
    Generalnego Gubernatorstwa w~1943~roku \cite{spis1943}.}
    \label{fig:spis_1943}
\end{figure}

W~marcu 1943~roku okupacyjne władze niemieckie dokonały spisu ludności
na terenie Generalnego Gubernatorstwa -- Desno było wówczas zakwalifikowane
przez niemieckich okupantów jako część gminy Okuniew w~dystrykcie
warszawskim. Spis ten wykazał, iż na terenie wsi Desno mieszkały wówczas
142~osoby. Wieś Desno była wówczas jedną z~najmniej licznych miejscowości
w~gminie Okuniew -- pod względem liczby mieszkańców przewyższała jedynie wsi:
Budziska (125~mieszkańców), Grabina (96~mieszkańców), Kazimierów
(115~mieszkańców), Konik Stary (66~mieszkańców), Królewskie Brzeziny
(79~mieszkańców) oraz Żwirówka (89~mieszkańców).

\begin{figure}[!ht]
    \vspace*{0.5cm}
    \centering \includegraphics[width=0.95\linewidth]{
        Desno_gmina_Halinow.png}
    \captionsetup{format=hang}
    \caption{Fragment rozporządzenia Prezesa Rady Ministrów Józefa
    Cyrankiewicza z~dnia 3~maja 1952~roku \cite{dziennik_ustaw2}.}
    \label{fig:halinow_1952}
\end{figure}

Pierwszego lipca 1952~roku, w~wyniku rozporządzenia Prezesa Rady Ministrów
Józefa Cyrankiewicza z~dnia 3~maja~1952~roku, wieś Desno, po blisko 57~latach
istnienia, wraz z~kilkoma okolicznymi wioskami, została wydzielona z~gminy
Okuniew i~przeniesiona do nowo powstałej gminy Halinów. Gmina Halinów nie
zaczęła wtedy funkcjonować jako samodzielna gmina, gdyż od razu po utworzeniu
została przekształcona w~dzielnicę powiatu miejsko-uzdrowiskowego Otwock
(patrz ryc. \ref{fig:halinow_1952}). W~następnych latach, w~wyniku kolejnych
reform samorządowych, wieś Desno najpierw stała się częścią gromady Halinów
w~powiecie otwockim województwa warszawskiego (w~latach 1958--1972) -- po
utworzeniu 1~stycznia 1958~roku powiatu otwockiego w~województwie warszawskim,
w~miejsce zlikwidowanego powiatu miejsko-uzdrowiskowego Otwock i~zamianie jego
dzielnic w~gromady i~osiedla. Następnie, na krótko, została częścią
reaktywowanej gminy Halinów w~ramach powiatu mińskiego województwa
warszawskiego (lata 1973--1974) po reformie z~1~stycznia 1973~roku likwidującej
gromady i~osiedla a~przywracającej w~ich miejsce gminy (patrz rys.
\ref{fig:desno_1972}).

\begin{figure}[!ht]
    \vspace*{0.5cm}
    \centering \includegraphics[width=1.0\linewidth]{
        Desno_1972.jpg}
    \captionsetup{format=hang}
    \caption{Fragment uchwały nr XX/93/72 Wojewódzkiej Rady Narodowej
    w~Warszawie z~dnia 1~grudnia 1972~roku \cite{dziennik_urzedowy2}.}
    \label{fig:desno_1972}
\end{figure}

Pierwszego czerwca 1975~roku wprowadzona została reforma likwidująca szczebel
powiatowy w~strukturach samorządowych w~Polskiej Rzeczpospolitej Ludowej, więc
gmina Halinów, a~wraz z~nią wieś Desno, zaczęła podlegać bezpośrednio pod
struktury nowo powstałego województwa stołecznego
warszawskiego\footnote{W~latach 1975--1998 wschodnia granica województwa
stołecznego warszawskiego przebiegała \emph{de facto} pomiędzy wsią Desno
a~wsią Olesin.}. Trwało to~przez ponad 23~lata, aż~do~końca 1998~roku, kiedy
to~reforma administracyjna III~Rzeczpospolitej, przywróciła wcześniej
zlikwidowany szczebel powiatowy oraz zmniejszyła liczbę województw
z~czterdziestu dziewięciu (wprowadzonych reformą z~1975~roku) do~szesnastu.
Finalnie, pierwszego stycznia 1999~roku, wieś Desno stała się częścią gminy
Halinów, w~powiecie mińskim, województwa mazowieckiego i~pozostaje w~tej
strukturze samorządowej aż~do~dnia pisania niniejszej książki.

\begin{figure}[!ht]
    \vspace*{0.5cm}
    \centering \includegraphics[width=1.0\linewidth]{
        parafia_sw_anny.jpg}
    \captionsetup{format=hang}
    \caption{Zdjęcie zabytkowego, XVI-wiecznego, drewnianego kościoła św. Anny
    w~Długiej Kościelnej, wykonane około 1963~roku \cite{galicka}.}
    \label{fig:sw_anna}
\end{figure}

Pierwszego marca 1987~roku kardynał Józef Glemp erygował parafię pod
wezwaniem Matki Łaski Bożej w~Halinowie i~tym samym wieś Desno, po blisko
92~latach istnienia, przestała być częścią parafii świętej Anny w~Długiej
Kościelnej -- została włączona do nowo powstałej parafii w~Halinowie. Ale nie na
długo, bo blisko ćwierć wieku później, 15~sierpnia 2011~roku, arcybiskup
Henryk Hoser erygował parafię pod wezwaniem Matki Odkupiciela w~Wielgolesie
Brzezińskim i~Desno stało się finalnie częścią tej właśnie parafii.

\newpage
\BigBold{Najważniejsze daty w~historii wsi \emph{Desna} / Desno:}

\begin{table}[h]
\centering
\begin{tabular}{>{\centering\arraybackslash}p{4cm} |
    >{\centering\arraybackslash}p{10cm}}
\textbf{ } & \textbf{ } \\ \hline
13~marca 1894~roku & Pierwsi mieszkańcy zaczynają sprowadzać się na teren
obecnego Desna (wtedy jeszcze części Cisia) \\
5~kwietnia 1894~roku & Założenie spółki \emph{Desna} przez 17~gospodarzy
chłopskich i~podpisanie umowy zakupu ziemi od Izraela Rotholca \\
11~lipca 1895~roku & Wydzielenie wsi \emph{Desna} z~dóbr ziemskich wsi Cisie \\
\textbf{23~lipca 1895~roku} & \textbf{Finalizacja zakupu wsi \emph{Desna} od
Izraela Rotholca przez współwłaścicieli spółki \emph{Desna}} \\
11~sieprnia 1895~roku & Pierwszy ślub we wsi \emph{Desna} --
Marianny Pauliny\indeksosoby[Piotrowska]{Marianna Paulina}{Kutrowska}{1875}{Wola
Rafałowska} Piotrowskiej oraz Franciszka
Kutrowskiego\indeksosoby{Franciszek}{Kutrowski}{1871}{Dąbrówka Grzybowska}\footnotemark \\
16~listopada 1895~roku & Narodziny pierwszej osoby we wsi \emph{Desna} --
Katarzyny Gomulskiej\indeksosoby[Gomulska]{Katarzyna}{Zawadzka}{1895}{Desno}
córki Stanisława i~Elżbiety \\
12~września 1896~roku & Zgon pierwszej osoby we wsi \emph{Desna} -- Pauliny
Piotrowskiej\indeksosoby{Paulina}{Piotrowska}{1896}{Desno}\\
29~lipca 1922~roku & Całkowita spłata pożyczki zaciągniętej przez
współwłaścicieli spółki \emph{Desna} \\
1~kwietnia 1930~roku & Wydzielenie wsi \emph{Desna} z~gminy Wiązowna
i~włączenie jej do gminy Okuniew \\
1~lipca 1952~roku & Wydzielenie wsi Desno z~gminy Okuniew i~włączenie jej do
gminy Halinów \\
1~marca 1987~roku & Włączenie wsi Desno do nowo powstałej parafii pw. Matki
Łaski Bożej w~Halinowie \\
15~sierpnia 2011~roku & Włączenie wsi Desno do nowo powstałej parafii pw. Matki
Odkupiciela w~Wielgolesie Brzezińskim \\
\end{tabular}
\end{table}

\footnotetext{W~akcie ślubu państwa Kutrowskich jest wprawdzie wpisane Cisie
jako miejsce ich zamieszkania, czyli sąsiednia wieś, z~której wieś
\emph{Desna} została wydzielona, ale wszystko wskazuje na to, iż oboje
mieszkali już wtedy na Deśnie -- musiało upłynąć trochę czasu zanim parafia
św.~Anny w~Długiej Kościelnej odnotowała formalne powstanie nowej wsi
i~zaczęła posługiwać się jej nazwą w~aktach metrykalnych.}

\newpage
\ifodd\value{page}\hbox{}\newpage\fi

\section*{Potomkowie Michała i~Marianny Gomulskich}

Pierwszym członkiem rodziny Gomulskich, który osiedlił się na terenie obecnego
Desna był Michał Gomulski\indeksosoby{Michał}{Gomulski}{1859}{Anielina}
(ur. 1859~r. -- zm. 1918~r.), przedostatni syn Piotra Franciszka Gomulskiego
i~Marianny z~domu Chłopik, młodszy brat Jana Gomulskiego, którego losy były
omawiane w~poprzednim podrozdziale ninijszej książki. Michał Gomulski
27~lutego 1881~roku ożenił się z~Marianną
Wróbel\indeksosoby[Wróbel]{Marianna}{Gomulska}{1860}{Siennica} (ur. 1860~r. --
zm. 1934~r.) w~kościele pod wezwaniem Narodzenia Najświętszej Maryi Panny
w~Mińsku Mazowieckim (patrz ryc. \ref{fig:mgomulski_1881}). Marianna Wróbel
urodziła się w~Siennicy (patrz ryc. \ref{fig:mwrobel_1860}), była córką
Jana Wróbla\indeksosoby{Jan}{Wróbel}{1815}{Siennica} (ur.~1815~r.~-- zm.~?),
pracującego jako służący w~tamtejszym probostwie, oraz Marianny
z~domu Drozd\indeksosoby[Drozd]{Marianna}{Wróbel}{1820}{Siennica}
(ur.~1820~r.~-- zm.~?). Marianna w~momencie ślubu z~Michałem mieszkała
w~Mikanowie, stąd też nowożeńcy przez pierwsze trzynaście lat małżeństwa
mieszkali właśnie w~tej miejscowości.

\begin{figure}[!ht]
    \vspace*{0.5cm}
    \centering \includegraphics[width=1.0\linewidth]{
        1860_Marianna_Wróbel_akt_chrztu_parafia_Siennica_wpis_90.jpg}
    \captionsetup{format=hang}
    \caption{Akt chrztu Marianny Wróbel -- par. Siennica 1860~rok (90/1860) 
    \cite{par_siennica}.}
    \label{fig:mwrobel_1860}
\end{figure}

Michał\indeksosoby{Michał}{Gomulski}{1859}{Anielina}
i~Marianna\indeksosoby[Wróbel]{Marianna}{Gomulska}{1860}{Siennica} Gomulscy
dochowali się jedenaściorga dzieci, z~czego pięcioro dożyło wieku
dorosłego\footnote{Drzewo genealogiczne najbliższej rodziny Michała i~Marianny
Gomulskich zostało zaprezentowane w~\hyperref[sec:drzewo_mgomulski]{załączniku
numer I}~do niniejszej książki.}:

\begin{itemize}
    \item Antoni Gomulski\indeksosoby{Antoni}{Gomulski}{1882}{Mikanów}
    (ur. 1882~r. -- zm. 1926~r.) -- patrz: ryc. \ref{fig:agomulski_1882},
    \item Stanisław Gomulski\indeksosoby{Stanisław}{Gomulski}{1883}{Mikanów}
    (ur. 1883~r. -- zm. 1884~r.) -- patrz: ryc. \ref{fig:sgomulski_1883},
    \item Marcin Gomulski\indeksosoby{Marcin}{Gomulski}{1885}{Mikanów}
    (ur. 1885~r. -- zm. ?) -- patrz: ryc. \ref{fig:mgomulski_1885},
    \item Jan Gomulski\indeksosoby{Jan}{Gomulski}{1888}{Mikanów}
    (ur. 1888~r. -- zm. 1923~r.) -- patrz: ryc. \ref{fig:jgomulski_1888},
    \item Władysław Gomulski\indeksosoby{Władysław}{Gomulski}{1890}{Mikanów}
    (ur. 1890~r. -- zm. 1890~r.) -- patrz: ryc. \ref{fig:wgomulski_1890},
    \item Stanisław Gomulski\indeksosoby{Stanisław}{Gomulski}{1892}{Mikanów}
    (ur. 1892~r. -- zm. 1961~r.) -- patrz: ryc. \ref{fig:sgomulski_1892},
    \item Zofia Gomulska\indeksosoby{Zofia}{Gomulska}{1894}{Cisie}
    (ur. 1894~r. -- zm. 1895~r.) -- patrz: ryc. \ref{fig:zgomulska_1894},
    \item Stefan Gomulski\indeksosoby{Stefan}{Gomulski}{1896}{Desno}
    (ur. 1896~r. -- zm. 1923~r.) -- patrz: ryc. \ref{fig:sgomulski_1896},
    \item Józef Gomulski\indeksosoby{Józef}{Gomulski}{1899}{Desno}
    (ur. 1899~r. -- zm. 1902~r.) -- patrz: ryc. \ref{fig:jgomulski_1899},
    \item Feliks Gomulski\indeksosoby{Feliks}{Gomulski}{1901}{Desno}
    (ur. 1901~r. -- zm. 1902~r.) -- patrz: ryc. \ref{fig:fgomulski_1901},
    \item Józefa Gomulska\indeksosoby[Gomulska]{Józefa}{Ładno}{1904}{Desno}[Lxyzadno]
    (ur. 1904~r. -- zm. 1987~r.) -- patrz: ryc. \ref{fig:jgomulska_1904}.
\end{itemize}

\begin{figure}[!ht]
    \vspace*{0.5cm}
    \centering \includegraphics[width=0.88\linewidth]{
        1881_Michał_Gumulski_Marianna_Wróbel_akt_ślubu_parafia_Mińsk_Mazowiecki_wpis_26.jpg}
    \captionsetup{format=hang}
    \caption{Akt ślubu Michała Gomulskiego oraz~Marianny Wróbel -- par. 
    Mińsk Mazowiecki 1881~rok (26/1881) \cite{par_minsk2}.}
    \label{fig:mgomulski_1881}
\end{figure}

\begin{figure}[!ht]
    \vspace*{0.5cm}
    \centering \includegraphics[width=0.77\linewidth]{
        1882_Antoni_Gomulski_akt_chrztu_parafia_Mińsk_Mazowiecki_wpis_21.jpg}
    \captionsetup{format=hang}
    \caption{Akt chrztu Antoniego Gomulskiego -- par. Mińsk Mazowiecki
    1882~rok (21/1882)
    \cite{par_minsk2}.}
    \label{fig:agomulski_1882}
\end{figure}

\begin{figure}[!ht]
    \vspace*{0.5cm}
    \centering \includegraphics[width=0.77\linewidth]{
        1883_Stanisław_Gomulski_akt_chrztu_parafia_Mińsk_Mazowiecki_wpis_264.jpg}
    \captionsetup{format=hang}
    \caption{Akt chrztu Stanisława Gomulskiego -- par. Mińsk Mazowiecki
    1883~rok (264/1883)
    \cite{par_minsk2}.}
    \label{fig:sgomulski_1883}
\end{figure}

\begin{figure}[!ht]
    \vspace*{0.5cm}
    \centering \includegraphics[width=0.77\linewidth]{
        1885_Marcin_Gomulski_akt_chrztu_parafia_Mińsk_Mazowiecki_wpis_300.jpg}
    \captionsetup{format=hang}
    \caption{Akt chrztu Marcina Gomulskiego -- par. Mińsk Mazowiecki
    1885~rok (300/1885)
    \cite{par_minsk2}.}
    \label{fig:mgomulski_1885}
\end{figure}

\begin{figure}[!ht]
    \vspace*{0.5cm}
    \centering \includegraphics[width=0.77\linewidth]{
        1888_Jan_Gomulski_akt_chrztu_parafia_Mińsk_Mazowiecki_wpis_144.jpg}
    \captionsetup{format=hang}
    \caption{Akt chrztu Jana Gomulskiego -- par. Mińsk Mazowiecki
    1888~rok (144/1888)
    \cite{par_minsk2}.}
    \label{fig:jgomulski_1888}
\end{figure}

\begin{figure}[!ht]
    \vspace*{0.5cm}
    \centering \includegraphics[width=0.72\linewidth]{
        1890_Władysław_Gomulski_akt_chrztu_parafia_Mińsk_Mazowiecki_wpis_315.jpg}
    \captionsetup{format=hang}
    \caption{Akt chrztu Władysława Gomulskiego -- par. Mińsk Mazowiecki
    1890~rok (315/1890)
    \cite{par_minsk2}.}
    \label{fig:wgomulski_1890}
\end{figure}

\begin{figure}[!ht]
    \vspace*{0.5cm}
    \centering \includegraphics[width=0.72\linewidth]{
        1892_Stanisław_Gomulski_akt_chrztu_parafia_Mińsk_Mazowiecki_wpis_47.jpg}
    \captionsetup{format=hang}
    \caption{Akt chrztu Stanisława Gomulskiego -- par. Mińsk Mazowiecki
    1892~rok (47/1892)
    \cite{par_minsk2}.}
    \label{fig:sgomulski_1892}
\end{figure}

\begin{figure}[!ht]
    \vspace*{0.5cm}
    \centering \includegraphics[width=0.77\linewidth]{
        1894_Zofia_Gomulska_akt_chrztu_parafia_Długa_Kościelna_wpis_118.jpg}
    \captionsetup{format=hang}
    \caption{Akt chrztu Zofii Gomulskiej -- par. Długa Kościelna
    1894~rok (118/1894)
    \cite{par_dluga}.}
    \label{fig:zgomulska_1894}
\end{figure}

\begin{figure}[!ht]
    \vspace*{0.5cm}
    \centering \includegraphics[width=0.77\linewidth]{
        1896_Stefan_Gomulski_akt_chrztu_parafia_Długa_Kościelna_wpis_56.jpg}
    \captionsetup{format=hang}
    \caption{Akt chrztu Stefena Gomulskiego -- par. Długa Kościelna
    1896~rok (56/1896)
    \cite{par_dluga}.}
    \label{fig:sgomulski_1896}
\end{figure}

\begin{figure}[!ht]
    \vspace*{0.5cm}
    \centering \includegraphics[width=0.82\linewidth]{
        1899_Józef_Gomulski_akt_chrztu_parafia_Długa_Kościelna_wpis_36.jpg}
    \captionsetup{format=hang}
    \caption{Akt chrztu Józefa Gomulskiego -- par. Długa Kościelna
    1899~rok (36/1899)
    \cite{par_dluga}.}
    \label{fig:jgomulski_1899}
\end{figure}

\begin{figure}[!ht]
    \vspace*{0.5cm}
    \centering \includegraphics[width=0.82\linewidth]{
        1901_Feliks_Gomulski_akt_chrztu_parafia_Długa_Kościelna_wpis_124.jpg}
    \captionsetup{format=hang}
    \caption{Akt chrztu Feliksa Gomulskiego -- par. Długa Kościelna
    1901~rok (124/1901)
    \cite{par_dluga}.}
    \label{fig:fgomulski_1901}
\end{figure}

\begin{figure}[!ht]
    \vspace*{0.5cm}
    \centering \includegraphics[width=0.87\linewidth]{
        1904_Józefa_Gomulska_Ładno_akt_chrztu_parafia_Długa_Kościelna_wpis_46.jpg}
    \captionsetup{format=hang}
    \caption{Akt chrztu Józefy Gomulskiej -- par. Długa Kościelna
    1904~rok (46/1904)
    \cite{par_dluga}.}
    \label{fig:jgomulska_1904}
\end{figure}

Około 1894~roku Michał\indeksosoby{Michał}{Gomulski}{1859}{Anielina}
i~Marianna\indeksosoby[Wróbel]{Marianna}{Gomulska}{1860}{Siennica} Gomulscy,
wraz z~czterema synami, przeprowadzili się do oddalonej o~blisko 18~kilometrów
na~północny zachód od~Mikanowa wsi Cisie (na~teren obecnego Desna), gdzie
przyszła na świat ich pierwsza córka
Zofia\indeksosoby{Zofia}{Gomulska}{1894}{Cisie} oraz czwórka ich najmłodszych
dzieci: Stefan\indeksosoby{Stefan}{Gomulski}{1896}{Desno},
Józef\indeksosoby{Józef}{Gomulski}{1899}{Desno},
Feliks\indeksosoby{Feliks}{Gomulski}{1901}{Desno} oraz
Józefa\indeksosoby[Gomulska]{Józefa}{Ładno}{1904}{Desno}[Lxyzadno] -- wszyscy
urodzeni już na~Deśnie. Spośród jedenaściorga dzieci Michała i~Marianny
Gomulskich, tylko sześcioro dożyło wieku dorosłego:
Antoni\indeksosoby{Antoni}{Gomulski}{1882}{Mikanów},
Marcin\indeksosoby{Marcin}{Gomulski}{1885}{Mikanów},
Jan\indeksosoby{Jan}{Gomulski}{1888}{Mikanów},
Stanisław\indeksosoby{Stanisław}{Gomulski}{1892}{Mikanów},
Stefan\indeksosoby{Stefan}{Gomulski}{1896}{Desno} oraz
Józefa\indeksosoby[Gomulska]{Józefa}{Ładno}{1904}{Desno}[Lxyzadno] , z~czego
tylko trójka osiedliła się na~stałe na~terenie obecnego Desna -- Antoni, Jan
oraz Stanisław, a~ich potomkowie mieszkają na~Deśnie do~dnia dzisiejszego.
\textbf{Antoni Gomulski}\indeksosoby{Antoni}{Gomulski}{1882}{Mikanów}
w~1921~roku ożenił się z~Magdaleną Zawadzką\footnote{Magdalena
Zawadzka\indeksosoby[Zawadzka][Gomulska]{Magdalena}{Poziemska}{1895}{Cisie}
była córką Marcina Zawadzkiego\indeksosoby{Marcin}{Zawadzki}{1856}{Anielina},
jednego ze~współudziałowców spółki \emph{Desna}, oraz prababką autora
niniejszej książki.} (ur. 1895~r. -- zm. 1970~r.) -- patrz: ryc.
\ref{fig:agomulski_1921}. Razem dochowali się jedynie dwójki synów
Romana\indeksosoby{Roman}{Gomulski}{1922}{Desno}\footnote{Roman Gomulski
w~1943~roku ożenił się ze~Stanisławą
Rogala\indeksosoby[Rogala]{Stanisława}{Gomulska}{1924}{Desno} (ur. 1924~r. --
zm. 2006~r.), córką Szymona
Rogali\indeksosoby{Szymon}{Rogala}{1878}{Kąty Goździejewskie}, jednego
ze~współudziałowców spółki \emph{Desna}, razem dochowali się piątki synów,
w~tym obecnego sołtysa wsi Desno -- Leszka Mariana
Gomulskiego\indeksosoby{Leszek Marian}{Gomulski}{1950}{Desno} (ur. 1950~r.).}
(ur.~1922~r.~-- zm.~1976~r.) oraz
Leona\indeksosoby{Leon}{Gomulski}{1924}{Desno} (ur.~1924~r.~-- zm. 1925~r.),
gdyż Antoni Gomulski\indeksosoby{Antoni}{Gomulski}{1882}{Mikanów} zmarł pięć
lat po~ich ślubie -- w~1926~roku. Magdalena
Zawadzka\indeksosoby[Zawadzka][Gomulska]{Magdalena}{Poziemska}{1895}{Cisie}
po~śmierci swojego pierwszego męża, w~1929~roku, wyszła ponownie za~mąż,
za~Michała
Poziemskiego\indeksosoby{Michał}{Poziemski}{1901}{Chobot}\footnote{Michał
Poziemski był synem Antoniego
Poziemskiego\indeksosoby{Antoni}{Poziemski}{1864}{Cisie},
jednego ze współudziałowców spółki \emph{Desna}, oraz pradziadkiem autora
niniejszej książki.} (ur. 1901~r. -- zm. 1947~r.) i~dochowali się wspólnie
czwórki dzieci:
Marianny\indeksosoby{Marianna}{Poziemska}{1931}{Desno} (ur.~1931~r.~--
zm.~1931~r.), Honoraty\indeksosoby{Honorata}{Poziemska}{1931}{Desno}
(ur.~1931~r.~-- zm.~1931~r.), Józefa\indeksosoby{Józefa}{Poziemska}{1933}{Desno}
(ur.~1933~r.~-- zm.~1933~r.) oraz Jadwigii
Leokadii\indeksosoby[Poziemska]{Jadwiga Leokadia}{Gomulska}{1935}{Desno}\footnote{Jadwiga
Leokadia Poziemska była babcią autora niniejszej książki.} (ur.~1935~r.~--
zm.~2014~r.).

\begin{figure}[!ht]
    \vspace*{0.5cm}
    \centering \includegraphics[width=0.95\linewidth]{
        1921_Antoni_Gomulski_Magdalena_Zawadzka_akt_ślubu_parafia_Długa_Kościelna_wpis_15.jpg}
    \captionsetup{format=hang}
    \caption{Akt ślubu Antoniego Gomulskiego oraz Magdaleny Zawadzkiej -- par.
    Długa Kościelna 1921~rok (15/1921)
    \cite{par_dluga}.}
    \label{fig:agomulski_1921}
\end{figure}

\textbf{Marcin Gomulski}\indeksosoby{Marcin}{Gomulski}{1885}{Mikanów}
w~1912~roku ożenił się z~Martą
Szulc\indeksosoby[Szulc]{Marta}{Gomulska}{1883}{Gilino} (ur.~1883~r.~--
zm.~?) i~zamieszkali razem w~sąsiędniej wsi -- w~Żwirówce, skąd pochodziła
Marta (patrz: ryc. \ref{fig:mgomulski_1912}). Dochowali się wspólnie
sześciorga dzieci:
Jerzego\indeksosoby{Jerzy}{Gomulski}{1912}{Żwirówka} (ur.~1912~r.~--
zm.~1988~r.), Wandy\indeksosoby{Wanda}{Gomulska}{1914}{Żwirówka}
(ur.~1914~r.~-- zm.~?), Karola
Henryka\indeksosoby{Karol Henryk}{Gomulski}{1915}{Żwirówka}
(ur.~1915~r.~-- zm.~?), Tadeusza
Stefana\indeksosoby{Tadeusz Stefan}{Gomulski}{1917}{Żwirówka} (ur.~1917~r.~--
zm.~1917~r.), Tadeusza\indeksosoby{Tadeusz}{Gomulski}{1921}{Żwirówka}
(ur.~1921~r.~-- zm.~2002~r.) oraz
Dominika\indeksosoby{Dominik}{Gomulski}{1925}{Żwirówka} (ur.~1925~r.~--
zm.~1925~r.).

\begin{figure}[!ht]
    \vspace*{0.5cm}
    \centering \includegraphics[width=0.95\linewidth]{
        1912_Marcin_Gomulski_Marta_Szulc_akt_ślubu_parafia_Długa_Kościelna_wpis_11.jpg}
    \captionsetup{format=hang}
    \caption{Akt ślubu Marcina Gomulskiego oraz Marty Szulc -- par.
    Długa Kościelna 1912~rok (11/1912)
    \cite{par_dluga}.}
    \label{fig:mgomulski_1912}
\end{figure}

\textbf{Jan Gomulski}\indeksosoby{Jan}{Gomulski}{1888}{Mikanów} w~1914~roku
ożenił się z~Rozalią Piekut\indeksosoby[Piekut]{Rozalia}{Gomulska}{1896}{Poręby}
(ur.~1896~r.~-- zm.~1934~r.) z~Poręb (patrz: ryc. \ref{fig:mgomulski_1912})
i~dochował się z~nią trójki córek:
Franciszki\indeksosoby[Gomulska]{Franciszka}{Gomulska}{1916}{Desno}
(ur.~1916~r.~-- zm.~1984~r.), Leokadii
Marianny\indeksosoby[Gomulska]{Leokodia Marianna}{Dobosz vel
Jerzak}{1918}{Desno} (ur.~1918~r.~-- zm.~?) oraz
Czesławy\indeksosoby[Gomulska]{Czesława}{Gomulska}{1922}{Desno} (ur.~1922~r.~--
zm.~2011~r.) -- wszystkie urodziły się na~Deśnie. W~1923~roku, zaledwie
dziewięć lat po~ślubie, Jan Gomulski\indeksosoby{Jan}{Gomulski}{1888}{Mikanów}
zmarł. Jego żona, Rozalia
Gomulska\indeksosoby[Piekut]{Rozalia}{Gomulska}{1896}{Poręby}, dwa lata
po~śmierci swojego pierwszego męża, wyszła ponownie za mąż --
za~Tomasza Gomulskiego\indeksosoby{Tomasz}{Gomulski}{1900}{Desno}
(ur.~1900~r.~-- zm.~1982~r.), odległego kuzyna Jana, syna
Stanisława\indeksosoby{Stanisław}{Gomulski}{1861}{Wólka Mińska}
i~Julianny\indeksosoby[Gągol]{Julianna}{Gomulska}{1862}{Dłużka} Gomulskich --
był to~pierwszy przypadek w~historii Desna i~prawdopodobnie pierwszy przypadek
w~historii rodziny Gomulskich (lecz nie ostatni), w~którym dwie osoby
o~nazwisku Gomulski wzięło ze~sobą ślub.

\begin{figure}[!ht]
    \vspace*{0.5cm}
    \centering \includegraphics[width=0.95\linewidth]{
        1914_Rozalia_Gomulska_Piekut_Jan_Gomulski_akt_ślubu_parafia_Pustelnik_wpis_22.jpg}
    \captionsetup{format=hang}
    \caption{Akt ślubu Jana Gomulskiego oraz Rozalii Piekut -- par. Pustelnik
    1914~rok (22/1914)
    \cite{par_pustelnik}.}
    \label{fig:jgomulski_1914}
\end{figure}

Szósty syn Michała i~Marianny Gomulskich,
\textbf{Stanisław}\indeksosoby{Stanisław}{Gomulski}{1892}{Mikanów},
w~1915~roku ożenił się z~Marianną
Kowalczyk\indeksosoby[Kowalczyk]{Marianna}{Gomulska}{1897}{Lubomin}
(ur.~1897~r.~-- zm.~1934~r.) z~Krzewiny (patrz: ryc. \ref{fig:sgomulski_1915}),
gdzie początkowo zamieszkali, by~około 1918~roku przeprowadzić się na~Desno.
Wspólnie dochowali się sześciorga dzieci:
Czesława\indeksosoby{Czesław}{Gomulski}{1915}{Krzewina} (ur.~1915~r.~--
zm.~1915~r.), Kazimiery\indeksosoby[Gomulska]{Kazimiera}{Popławska}{1916}{Krzewina}
(ur.~1916~r.~-- zm.~?), Czesława\indeksosoby{Czesław}{Gomulski}{1918}{Desno}
(ur.~1918~r.~-- zm.~1918~r.), Heleny\indeksosoby[Gomulska]{Helena}{Komorowska}{1921}{Desno}
(ur.~1921~r.~-- zm.~?), Janiny\indeksosoby{Janina}{Gomulski}{1923}{Desno}
(ur.~1923~r.~-- zm.~1924~r.) oraz
Henryka\indeksosoby{Henryk}{Gomulski}{1927}{Desno} (ur.~1927~r.~--
zm.~1993~r.).

\begin{figure}[!ht]
    \vspace*{0.5cm}
    \centering \includegraphics[width=0.8\linewidth]{
        1915_Stanisław_Gomulski_Marianna_Kowalczyk_akt_ślubu_parafia_Długa_Kościelna_wpis_2.jpg}
    \captionsetup{format=hang}
    \caption{Akt ślubu Stanisława Gomulskiego oraz Marianny Kowalczyk -- par.
    Długa Kościelna 1915~rok (2/1915)
    \cite{par_dluga}.}
    \label{fig:sgomulski_1915}
\end{figure}

Siódmy syn Michała i~Marianny Gomulskich,
\textbf{Stefan}\indeksosoby{Stefan}{Gomulski}{1896}{Desno}, w~1918~roku
ożenił się, ze~starszą od~niego o~osiem lat, Magdaleną
Tomaszewską\indeksosoby[Jackiewicz][Tomaszewska]{Magdalena}{Gomulska}{1888}{Desno}
z~domu Jackiewicz (ur.~1888~r.~-- zm.~1952~r.) wdową po~Tomaszu
Tomaszewskim\indeksosoby{Tomasz}{Tomaszewski}{1883}{Chobot} (ur.~1883~r.~--
zm.~1913~r.) z~Chobotu (patrz: ryc. \ref{fig:sgomulski_1918}). Para doczekała
się córki Stefanii\indeksosoby{Stefania}{Gomulska}{1920}{Chobot}
(ur.~1920~r.~-- zm.~?) oraz syna
Stefana\indeksosoby{Stefan}{Gomulski}{1923}{Chobot} (ur.~1923~r.~--
zm.~1998~r.), który urodził się w~roku śmierci swojego ojca. Stefan Gomulski
(senior) zmarł w~1923~roku, zaledwie pięć lat po~swoim ślubie.

\begin{figure}[!ht]
    \vspace*{0.5cm}
    \centering \includegraphics[width=0.9\linewidth]{
        1918_Stefan_Gomulski_Magdalena_Tomaszewska_akt_ślubu_parafia_Długa_Kościelna_wpis_3.jpg}
    \captionsetup{format=hang}
    \caption{Akt ślubu Stefana Gomulskiego oraz Magdaleny Tomaszewskiej -- par.
    Długa Kościelna 1918~rok (3/1918)
    \cite{par_dluga}.}
    \label{fig:sgomulski_1918}
\end{figure}

Stefan Gomulski\indeksosoby{Stefan}{Gomulski}{1896}{Desno} był
najprawdopodobniej ojcem nieślubnego dziecka Magdaleny
Zawadzkiej\indeksosoby[Zawadzka][Gomulska]{Magdalena}{Poziemska}{1895}{Cisie}
-- Marjana\indeksosoby{Marjan}{Zawadzki}{1917}{Desno} (ur.~1917~r.~-- zm.~?).
Dnia 9~września 1917~roku we wsi \emph{Desna} na~świat przyszedł Marjan
Zawadzki, faktu jego urodzenia, nie zgłosił księdzu parafii w~Długiej
Kościelnej jego ojciec, jak to było w~zwyczaju, ale przyszły ojciec chrzestny
-- Franciszek Gomulski\indeksosoby{Franciszek}{Gomulski}{1897}{Desno} (syn
Stanisława i~Julianny Gomulskich), daleki kuzyn omawianego tu~Stefana
Gomulskiego. W~księgach parafialnych Marjan został zarejestrowany pod
nazwiskiem matki (patrz: ryc. \ref{fig:marzawadzki_1917}), co było standardową
praktyką w~tamtych czasach w~przypadku dzieci urodzonych poza małżeństwem.

\begin{figure}[!ht]
    \vspace*{0.5cm}
    \centering \includegraphics[width=0.97\linewidth]{
        1917_Marian_Zawadzki_akt_chrztu_parafia_Długa_Kościelna_wpis_60.jpg}
    \captionsetup{format=hang}
    \caption{Akt chrztu Marjana Zawadzkiego -- par. Długa Kościelna 1917~rok
    (60/1917)
    \cite{par_dluga}.}
    \label{fig:marzawadzki_1917}
\end{figure}

Dnia 11~października 1917~roku, nieco ponad miesiąc po~urodzeniu
Marjana\indeksosoby{Marjan}{Zawadzki}{1917}{Desno},
Magdalena
Zawadzka\indeksosoby[Zawadzka][Gomulska]{Magdalena}{Poziemska}{1895}{Cisie}
złożyła w~sądzie gminnym w~Otwocku powództwo o~alimenty dla
dziecka, wskazując w~nim Stefana
Gomulskiego\indeksosoby{Stefan}{Gomulski}{1896}{Desno} jako ojca Marjana
Zawadzkiego (patrz: ryc. \ref{fig:powodztwo_1917}). Sąd Pokoju w~Otwocku,
po~przeprowadzeniu postępowania dowodowego, dnia 12~lutego 1918~roku,
przychylił się do~powództwa Magdaleny Zawadzkiej i~nakazał Stefanowi
Gomulskiemu zapłacić 1000~marek alimentów na rzecz ich wspólnego dziecka
(patrz: ryc. \ref{fig:wyrok_1918}). Wyrok był zaoczny -- Stefan Gomulski nie
stawił się na~rozprawie, a~o~wyroku miał się dowiedzieć dopiero, gdy komornik
zajął mu~konia. Powództwo Magdaleny Zawadzkiej zostało finalnie oddalone przez
wyrok Sądu Pokoju z~dnia 25~października 1918~roku, po~tym gdy Stefan Gomulski
wynajął adwokata, który podważył wiarygodność zeznań świadków. Mniej niż trzy
lata później Magdalena Zawadzka wyszła za~mąż
za~Antoniego\indeksosoby{Antoni}{Gomulski}{1882}{Mikanów} Gomulskiego
(najstarszego brata Stefana) i~wspólnie wychowywali Marjana razem
z~pozostałymi swoimi dziećmi.

\begin{figure}[!ht]
    \vspace*{0.5cm}
    \centering \includegraphics[width=0.9\linewidth]{
        Magdalena_Zawadzka_powodztwo_o_alimenty.jpg}
    \captionsetup{format=hang}
    \caption{Powództwo Magdaleny Zawadzkiej przeciwko Stefanowi Gomulskiemu
    o~alimenty na rzecz Marjana Zawadzkiego \cite{otwock}.}
    \label{fig:powodztwo_1917}
\end{figure}

\begin{figure}[!ht]
    \vspace*{0.5cm}
    \centering \includegraphics[width=0.93\linewidth]{
        Marjan_Zawadzki_alimenty_wyrok.jpg}
    \captionsetup{format=hang}
    \caption{Wyrok Sądu Pokoju w~Otwocku z~dnia 12~lutego 1918~roku
    \cite{otwock}.}
    \label{fig:wyrok_1918}
\end{figure}

Ostatnim dzieckiem Michała i~Marianny Gomulskich, które dożyło wieku dorosłego,
była \textbf{Józefa}\indeksosoby[Gomulska]{Józefa}{Ładno}{1904}{Desno}[Lxyzadno],
która w~1921~roku wyszła za~mąż za~Franciszka
Ładno\indeksosoby{Franciszek}{Ładno}{1898}{Cisie}[Lxyzadno] (ur.~1898~r.~--
zm.~1961~r.) z~Cisia (patrz: ryc. \ref{fig:jgomulska_1921}) i~tam też razem
zamieszkali po~ślubie. Dochowali się wspólnie prawdopodobnie tylko dwóch córek:
Leokadii\indeksosoby{Leokadia}{Ładno}{1923}{Cisie}[Lxyzadno] (ur.~1923~r.~--
zm.~1994~r.) oraz Zofii\indeksosoby{Zofia}{Ładno}{1926}{Cisie}[Lxyzadno]
(ur.~1926~r.~-- zm.~2007~r.).

\begin{figure}[!ht]
    \vspace*{0.5cm}
    \centering \includegraphics[width=0.87\linewidth]{
        1921_Józefa_Gomulska_Franciszek_Ładno_akt_ślubu_parafia_Długa_Kościelna_wpis_28.jpg}
    \captionsetup{format=hang}
    \caption{Akt ślubu Józefy Gomulskiej i~Franciszka Ładno -- par. Długa
    Kościelna 1921~rok (28/1921)
    \cite{par_dluga}.}
    \label{fig:jgomulska_1921}
\end{figure}

Michał Gomulski\indeksosoby{Michał}{Gomulski}{1859}{Anielina} zmarł
11~sierpnia 1918~roku w~wieku niespełna 59~lat (patrz: ryc.
\ref{fig:mgomulski_1918}), natomiast jego żona
Marianna\indeksosoby[Wróbel]{Marianna}{Gomulska}{1860}{Siennica} żyła jeszcze
przez ponad 16~lat po~jego śmierci,umierając 29~października 1934~roku,
w~wieku ponad 74~lat (patrz: ryc. \ref{fig:mgomulska_1934}). Oboje zostali
prawdopodobnie pochowani na~cmentarzu parafialnym parafii św.~Anny w~Długiej
Kościelnej, ale nie sposób jest odnaleźć ich grobów na~aktualnym terenie
cmentarza. Większość potomków rodziny Michała i~Marianny Gomulskich spoczywa
prawdopodobnie na~terenie cmentarza w~Długiej Kościelnej, gdyż ich dzieci nie
wyprowadziły się zbyt daleko od~rodzinnej wsi. Ryciny: \ref{fig:rgomulski_2006},
\ref{fig:rgomulska_2008}, \ref{fig:hgomulski_2003}, \ref{fig:sgomulski_1923},
\ref{fig:jladno_2007} oraz \ref{fig:pgomulski_2019} przedstawiają przykłady
niektórych z~ich nagrobków.

\begin{figure}[!ht]
    \vspace*{0.5cm}
    \centering \includegraphics[width=0.83\linewidth]{
        1918_Michał_Gomulski_akt_zgonu_parafia_Długa_Kościelna_wpis_34.jpg}
    \captionsetup{format=hang}
    \caption{Akt zgonu Michała Gomulskiego -- par. Długa Kościelna 1918~rok
    (34/1918) \cite{par_dluga}.}
    \label{fig:mgomulski_1918}
\end{figure}

\begin{figure}[!ht]
    \vspace*{0.5cm}
    \centering \includegraphics[width=0.83\linewidth]{
        1934_Marianna_Gumulska_Wróbel_akt_zgonu_parafia_Długa_Kościelna_wpis_43.jpg}
    \captionsetup{format=hang}
    \caption{Akt zgonu Marianny Gomulskiej -- par. Długa Kościelna 1934~rok
    (43/1934) \cite{par_dluga}.}
    \label{fig:mgomulska_1934}
\end{figure}

\begin{figure}[!ht]
    \vspace*{0.4cm}
    \centering \includegraphics[width=0.75\linewidth]{
        2006_Roman_Gomulski_Stanisława_Gomulska_zdjęcie_nagrobka2_cmentarz_Długa_Kościelna.jpg}
    \captionsetup{format=hang}
    \caption{Zdjęcie grobu małżonków Romana oraz Stanisławy Gomulskich,
    znajdującego się na cmentarzu parafialnym parafii św. Anny w~Długiej
    Kościelnej (sektor: A5, rząd: 10, numer: 14)}
    \label{fig:rgomulski_2006}
\end{figure}

\begin{figure}[!ht]
    \vspace*{0.4cm}
    \centering \includegraphics[width=0.79\linewidth]{
        2008_Tomasz_Gomulski_Franciszka_Gomulska_zdjęcie_nagrobka1_cmentarz_Długa_Kościelna.jpg}
    \captionsetup{format=hang}
    \caption{Zdjęcie grobu Rozalii Gomulskiej oraz członków jej najbliższej
    rodziny, znajdującego się na cmentarzu parafialnym parafii św. Anny
    w~Długiej Kościelnej (sektor: A5, rząd: 2, numer: 12)}
    \label{fig:rgomulska_2008}
\end{figure}

\begin{figure}[!ht]
    \vspace*{0.4cm}
    \centering \includegraphics[width=0.75\linewidth]{
        2003_Henryk_Gomulski_Cecylia_Gomulska_zdjęcie_nagrobka2_cmentarz_Długa_Kościelna.jpg}
    \captionsetup{format=hang}
    \caption{Zdjęcie grobu małżonków Henryka oraz Cecyli Gomulskich
    znajdującego, się na cmentarzu parafialnym parafii św. Anny w~Długiej
    Kościelnej (sektor: A5, rząd: 3, numer: 10)}
    \label{fig:hgomulski_2003}
\end{figure}

\begin{figure}[!ht]
    \vspace*{0.4cm}
    \centering \includegraphics[width=0.78\linewidth]{
        1923_Stefan_Gomulski_zdjęcie_nagrobka3_cmentarz_Długa_Kościelna.jpg}
    \captionsetup{format=hang}
    \caption{Zdjęcie grobu małżonków Stefana oraz Magdaleny Gomulskich,
    znajdującego się na cmentarzu parafialnym parafii św. Anny w~Długiej
    Kościelnej (sektor: B1, rząd: 14, numer: 7)}
    \label{fig:sgomulski_1923}
\end{figure}

\begin{figure}[!ht]
    \vspace*{0.4cm}
    \centering \includegraphics[width=0.78\linewidth]{
        2007_Józefa_Ładno_Gomulska_zdjęcie_nagrobka_cmentarz_Długa_Kościelna.jpg}
    \captionsetup{format=hang}
    \caption{Zdjęcie grobu małżonków Franciszka oraz~Józefy Ładno (z~domu
    Gomulskiej), znajdującego się na~cmentarzu parafialnym parafii św.~Anny
    w~Długiej Kościelnej (sektor:~A3, rząd:~6, numer:~19)}
    \label{fig:jladno_2007}
\end{figure}

\begin{figure}[!ht]
    \vspace*{0.4cm}
    \centering \includegraphics[width=0.78\linewidth]{
        2019_Piotr_Gomulski_Tadeusz_Gomulski_zdjęcie_nagrobka1_cmentarz_Długa_Kościelna.jpg}
    \captionsetup{format=hang}
    \caption{Zdjęcie grobu Tadeusza oraz~Piotra Gomulskich, znajdującego się
    na~cmentarzu parafialnym parafii św.~Anny w~Długiej Kościelnej (sektor:~A6,
    rząd:~12, numer:~32)}
    \label{fig:pgomulski_2019}
\end{figure}

\begin{figure}[!ht]
\vspace*{0.5cm}
\centering
\begin{tikzpicture}[
  % Definiujemy stałe odstępy
  node distance=1cm and 1cm, 
  every node/.style={draw, rectangle, align=center, minimum width=3.8cm,
  inner sep=5pt}]

% --- POKOLENIE 0 ---
\node (Osiowa0) [draw=none, minimum width=0pt, inner sep=0pt] at (0,0) {};
\node (Antoni) [left=1.5cm of Osiowa0] {Antoni Gomoła \\ {
    \footnotesize \emph{Wnętrzne, 1764 r.}}};
\node (Petronella) [right=1.5cm of Osiowa0] {Petronella Sokół \\ {
    \footnotesize \emph{Mysłów, 1768 r.}}};

\draw[-] (Antoni) -- (Petronella)
    node[midway, below, draw=none, minimum width=0pt,
    font=\scriptsize] (Slub0) {\emph{Wilczyska, 1783 r.}};

% --- POKOLENIE 1 ---
\node (Osiowa1) [draw=none, minimum width=0pt, inner sep=0pt,
below=1.5cm of Slub0] {};
\node (Jan) [left=1.5cm of Osiowa1] {Jan Gomoła \\ {
    \footnotesize \emph{Mysłów, 1789 r.}}};
\node (Łucja) [right=1.5cm of Osiowa1] {Łucja Rokicka \\ {
    \footnotesize \emph{Rudzienko, 1790 r.}}};

\draw[-] (Jan) -- (Łucja)
    node[midway, above, draw=none, minimum width=0pt,
    font=\scriptsize] (Slub1_Gora) { }
    node[midway, below, draw=none, minimum width=0pt,
    font=\scriptsize] (Slub1) {\emph{Kołbiel, 1814 r.}};

\draw[->] (Slub0) -- (Slub1_Gora);

% --- POKOLENIE 2 ---
\node (Osiowa2) [draw=none, minimum width=0pt, inner sep=0pt,
below=1.5cm of Slub1] {};
\node (Piotr) [left=1.5cm of Osiowa2] {Piotr Franciszek Gumulski \\ {
    \footnotesize \emph{Rudzienko, 1815 r.}}};
\node (Marianna) [right=1.5cm of Osiowa2] {Marianna Chłopik \\ {
    \footnotesize \emph{Zamienie, 1819 r.}}};
    
\draw[-] (Piotr) -- (Marianna)
    node[midway, above, draw=none, minimum width=0pt,
    font=\scriptsize] (Slub2_Gora) { }
    node[midway, below, draw=none, minimum width=0pt,
    font=\scriptsize] (Slub2) {\emph{Mińsk, 1835 r.}};

\draw[->] (Slub1) -- (Slub2_Gora);

% --- POKOLENIE 3 ---
\node (Osiowa3) [draw=none, minimum width=0pt, inner sep=0pt,
below=1.5cm of  Slub2] {};
\node (Michał) [left=1.5cm of Osiowa3] {Michał Gomulski \\ {
    \footnotesize \emph{Anielina, 1859 r.}}};
\node (Marianna2) [right=1.5cm of Osiowa3] {Marianna Wróbel \\ {
    \footnotesize \emph{Siennica, 1860 r.}}};

\draw[-] (Michał) -- (Marianna2)
    node[midway, above, draw=none, minimum width=0pt,
    font=\scriptsize] (Slub3_Gora) { }
    node[midway, below, draw=none, minimum width=0pt,
    font=\scriptsize] (Slub3) {\emph{Mińsk, 1881 r.}};

\draw[->] (Slub2) -- (Slub3_Gora);

\end{tikzpicture}
\end{figure}

Bezpośredni potomkowie Michała i~Marianny Gomulskich, mieszkają na terenie
obecnego Desna do dnia dzisiejszego -- większość z~nich nosi nazwisko Gomulski,
ale poza nimi, Desno zamieszkują też dwie rodziny o~nazwisku Bešić, które
również wywodzą się z~tej, najstarszej na Deśnie, gałęzi rodziny Gomulskich.

\newpage
\ifodd\value{page}\hbox{}\newpage\fi

\section*{Potomkowie Stanisława i~Elżbiety Gomulskich}

Drugim członkiem rodziny Gomulskich, który osiedlił się na terenie obecnego
Desna, był Stanisław Gomulski\indeksosoby{Stanisław}{Gomulski}{1863}{Anielina}
(ur. 1863~r. -- zm. 1929~r.), najmłodszy syn Piotra Franciszka Gomulskiego
i~Marianny z~domu Chłopik, młodszy brat Jana i~Michała. Stanisław Gomulski
23~stycznia 1882~roku ożenił się z~Elżbietą
Brodzińską\indeksosoby[Brodzińska]{Elżbieta}{Gomulska}{1857}{Grębiszew}
(ur. 1857~r. -- zm. 1928~r.) w~kościele pod wezwaniem Narodzenia Najświętszej
Maryi Panny w~Mińsku Mazowieckim (patrz ryc. \ref{fig:sgomulski_1882}).
Elżbieta Brodzińska urodziła się w~1857~roku w~Grębiszewie (parafia Glinianka),
była córką Grzegorza
Brodzińskiego\indeksosoby{Grzegorz}{Brodziński}{1833}{Rudzienko}
(ur.~1833~r.~-- zm. 1862~r.) oraz
Marianny\indeksosoby[Mazek]{Marianna}{Brodzińska}{1836}{Podgórzno} z~domu
Mazek (ur. 1836~r. -- zm. 1881~r.) -- patrz ryc. \ref{fig:ebrodzinska_1857}.
Para młoda przez pierwsze kilka lat po~ślubie mieszkała we~wsi Jarząbek, gdzie
urodziła się trójka ich najstarszych dzieci:
Jan\indeksosoby{Jan}{Gomulski}{1882}{Jarząbek},
Marianna\indeksosoby[Gomulska]{Marianna}{Świtkowska}{1884}{Jarząbek}[Sxyzwitkowska]
oraz Antonina\indeksosoby[Gomulska]{Antonina}{Foryś}{1887}{Jarząbek}, żeby
około 1890~roku powrócić do~rodzinnej wsi Stanisława -- do Anielny, gdzie
urodziła się kolejna dwójka ich dzieci
(Stanisława\indeksosoby[Gomulska]{Stanisława}{Bronikowska}{1890}{Anielina}
i~Józef\indeksosoby{Józef}{Gomulski}{1892}{Anielina}).

\begin{figure}[!ht]
    \vspace*{0.5cm}
    \centering \includegraphics[width=0.9\linewidth]{
        1857_Elżbieta_Gumulska_Brodzińska_akt_chrztu_parafia_Glinianka_wpis_55.jpg}
    \captionsetup{format=hang}
    \caption{Akt chrztu Elżbiety Brodzińskiej -- par. Glinianka 1857~rok (55/1857)
    \cite{par_glinianka}.}
    \label{fig:ebrodzinska_1857}
\end{figure}

Stanisław i~Elżbieta Gomulscy dochowali się ośmiorga dzieci, z~czego aż
siedmioro dożyło wieku dorosłego\footnote{Drzewo genealogiczne najbliższej
rodziny Stanisława i~Elżbiety Gomulskich zostało zaprezentowane
w~\hyperref[sec:drzewo_sgomulski1]{załączniku numer I}~do niniejszej książki.}:

\begin{itemize}
    \item Jan Gomulski\indeksosoby{Jan}{Gomulski}{1882}{Jarząbek}
    (ur. 1882~r. -- zm. 1943~r.) -- patrz: ryc. \ref{fig:jgomulski_1882},
    \item Marianna Gomulska\indeksosoby[Gomulska]{Marianna}{Świtkowska}{1884}{Jarząbek}[Sxyzwitkowska]
    (ur. 1884~r. -- zm. 1954~r.) -- patrz: ryc. \ref{fig:mgomulska_1884},
    \item Antonina Gomulska\indeksosoby[Gomulska]{Antonina}{Foryś}{1887}{Jarząbek}
    (ur. 1887~r. -- zm. 1971) -- patrz: ryc. \ref{fig:agomulska_1887},
    \item Stanisława Gomulska\indeksosoby[Gomulska]{Stanisława}{Bronikowska}{1890}{Anielina}
    (ur. 1890~r. -- zm. 1976~r.) -- patrz: ryc. \ref{fig:sgomulska_1890},
    \item Józef Gomulski\indeksosoby{Józef}{Gomulski}{1892}{Anielina}
    (ur. 1892~r. -- zm. 1929~r.) -- patrz: ryc. \ref{fig:jgomulski_1892},
    \item Katarzyna Gomulska\indeksosoby[Gomulska]{Katarzyna}{Zawadzka}{1895}{Desno}
    (ur. 1895~r. -- zm. 1987~r.) -- patrz: ryc. \ref{fig:kgomulska_1895},
    \item Andrzej Gomulski\indeksosoby{Andrzej}{Gomulski}{1898}{Desno}
    (ur. 1898~r. -- zm. 1902~r.) -- patrz: ryc. \ref{fig:agomulski_1898},
    \item Władysław Gomulski\indeksosoby{Władysław}{Gomulski}{1901}{Desno}
    (ur. 1901~r. -- zm. 1967~r.) -- patrz: ryc. \ref{fig:wgomulski_1901},
\end{itemize}

\begin{figure}[!ht]
    \vspace*{0.5cm}
    \centering \includegraphics[width=0.85\linewidth]{
        1882_Stanisław_Gumulski_Elżbieta_Brodzińska_akt_ślubu_parafia_Mińsk_Mazowiecki_wpis_15.jpg}
    \captionsetup{format=hang}
    \caption{Akt ślubu Stanisława Gomulskiego oraz Elżbiety Brodzińskiej --
    par. Mińsk Mazowiecki 1882~rok (15/1882) \cite{par_minsk2}.}
    \label{fig:sgomulski_1882}
\end{figure}

\begin{figure}[!ht]
    \vspace*{0.5cm}
    \centering \includegraphics[width=0.77\linewidth]{
        1882_Jan_Gomulski_akt_chrztu_parafia_Mińsk_Mazowiecki_wpis_360.jpg}
    \captionsetup{format=hang}
    \caption{Akt chrztu Jana Gomulskiego -- par. Mińsk Mazowiecki
    1882~rok (360/1882)
    \cite{par_minsk2}.}
    \label{fig:jgomulski_1882}
\end{figure}

\begin{figure}[!ht]
    \vspace*{0.5cm}
    \centering \includegraphics[width=0.77\linewidth]{
        1884_Marianna_Gomulska_Świtkowska_akt_chrztu_parafia_Mińsk_Mazowiecki_wpis_377.jpg}
    \captionsetup{format=hang}
    \caption{Akt chrztu Marianny Gomulskiej -- par. Mińsk Mazowiecki
    1884~rok (377/1884)
    \cite{par_minsk2}.}
    \label{fig:mgomulska_1884}
\end{figure}

\begin{figure}[!ht]
    \vspace*{0.5cm}
    \centering \includegraphics[width=0.75\linewidth]{
        1887_Antonina_Gomulska_Foryś_akt_chrztu_parafia_Mińsk_Mazowiecki_wpis_226.jpg}
    \captionsetup{format=hang}
    \caption{Akt chrztu Antoniny Gomulskiej -- par. Mińsk Mazowiecki
    1887~rok (226/1887)
    \cite{par_minsk2}.}
    \label{fig:agomulska_1887}
\end{figure}

\begin{figure}[!ht]
    \vspace*{0.5cm}
    \centering \includegraphics[width=0.75\linewidth]{
        1890_Stanisława_Gomulska_Bronikowska_akt_chrztu_parafia_Mińsk_Mazowiecki_wpis_19.jpg}
    \captionsetup{format=hang}
    \caption{Akt chrztu Stanisławy Gomulskiej -- par. Mińsk Mazowiecki
    1890~rok (19/1890)
    \cite{par_minsk2}.}
    \label{fig:sgomulska_1890}
\end{figure}

\begin{figure}[!ht]
    \vspace*{0.5cm}
    \centering \includegraphics[width=0.75\linewidth]{
        1892_Józef_Gomulski_akt_chrztu_parafia_Mińsk_Mazowiecki_wpis_408.jpg}
    \captionsetup{format=hang}
    \caption{Akt chrztu Józefa Gomulskiego -- par. Mińsk Mazowiecki
    1892~rok (408/1892)
    \cite{par_minsk2}.}
    \label{fig:jgomulski_1892}
\end{figure}

\begin{figure}[!ht]
    \vspace*{0.5cm}
    \centering \includegraphics[width=0.75\linewidth]{
        1895_Katarzyna_Gomulska_Zawadzka_akt_chrztu_parafia_Długa_Kościelna_wpis_108.jpg}
    \captionsetup{format=hang}
    \caption{Akt chrztu Katarzyny Gomulskiej -- par. Długa Kościelna
    1895~rok (108/1895)
    \cite{par_dluga}.}
    \label{fig:kgomulska_1895}
\end{figure}

\begin{figure}[!ht]
    \vspace*{0.5cm}
    \centering \includegraphics[width=0.75\linewidth]{
        1898_Andrzej_Gomulski_akt_chrztu_parafia_Długa_Kościelna_wpis_93.jpg}
    \captionsetup{format=hang}
    \caption{Akt chrztu Andrzeja Gomulskiego -- par. Długa Kościelna
    1898~rok (93/1898)
    \cite{par_dluga}.}
    \label{fig:agomulski_1898}
\end{figure}

\begin{figure}[!ht]
    \vspace*{0.5cm}
    \centering \includegraphics[width=0.75\linewidth]{
        1901_Władysław_Gomulski_akt_chrztu_parafia_Długa_Kościelna_wpis_49.jpg}
    \captionsetup{format=hang}
    \caption{Akt chrztu Władysłwa Gomulskiego -- par. Długa Kościelna
    1901~rok (49/1901)
    \cite{par_dluga}.}
    \label{fig:wgomulski_1901}
\end{figure}

Około 1895~roku
Stanisław\indeksosoby{Stanisław}{Gomulski}{1863}{Anielina}
i~Elżbieta\indeksosoby[Brodzińska]{Elżbieta}{Gomulska}{1857}{Grębiszew}
Gomulscy, wraz z~piątką dzieci, przeprowadzili się do~oddalonej o~blisko
16~kilometrów na północny zachód od~Anieliny wsi \emph{Desna}, gdzie przyszła
na~świat trójka ich najmłodszych dzieci:
Katarzyna\indeksosoby[Gomulska]{Katarzyna}{Zawadzka}{1895}{Desno},
Andrzej\indeksosoby{Andrzej}{Gomulski}{1898}{Desno}, oraz
Władysław\indeksosoby{Władysław}{Gomulski}{1901}{Desno}. Spośród ich ośmiorga
dzieci, tylko dwóch synów osiedliło się na stałe na terenie obecnego Desna --
najstarszy syn
Jan\indeksosoby{Jan}{Gomulski}{1882}{Jarząbek} oraz jego młodszy brat
Józef\indeksosoby{Józef}{Gomulski}{1892}{Anielina}. \textbf{Jan Gomulski}
w~1912~roku ożenił się z~Eleonorą
Kosieradzką\indeksosoby[Kosieradzka]{Eleonora}{Gomulska}{1894}{Warszawa}
(ur. 1894~r. -- zm. 1972~r.), osiemnastolatką z~sąsiedniej wsi Krzewina,
córką Teofila\indeksosoby{Teofil}{Kosieradzki}{1865}{Warszawa}
(ur.~1865~r.~-- zm.~1917~r.) i~Wiktorii
Ewy\indeksosoby[Kozłowska]{Wiktoria Ewa}{Kosieradzka}{1870}{Warszawa} z~domu
Kozłowskiej (ur.~1870~r.~-- zm.~1949~r.) Kosieradzkich. (patrz ryc.
\ref{fig:jgomulski_1912}).

\begin{figure}[!ht]
    \vspace*{0.5cm}
    \centering \includegraphics[width=0.9\linewidth]{
        1912_Jan_Gomulski_Eleonora_Kosieradzka_akt_ślubu_parafia_Długa_Kościelna_wpis_2.jpg}
    \captionsetup{format=hang}
    \caption{Akt ślubu Jana Gomulskiego oraz Eleonory Kosieradzkiej -- par. 
    Długa Kościelna 1912~rok (2/1912) \cite{par_dluga}.}
    \label{fig:jgomulski_1912}
\end{figure}

Jan\indeksosoby{Jan}{Gomulski}{1882}{Jarząbek}
i~Eleonora\indeksosoby[Kosieradzka]{Eleonora}{Gomulska}{1894}{Warszawa}
Gomulscy dochowali się razem sześciorga dzieci: Natalii
Zofii\indeksosoby{Natalia Zofia}{Gomulska}{1912}{Desno}
(ur.~1912~r.~-- zm.~?), Wiktorii\indeksosoby[Gomulska]{Wiktoria}{Zawada}{1916}{Desno}
(ur.~1916~r.~-- zm.~2009~r.), Ireny\indeksosoby{Irena}{Gomulska}{1919}{Desno}
(ur.~1919~r.~-- zm. 1920~r.), Henryka\indeksosoby{Henryk}{Gomulski}{1922}{Desno}
(ur.~1922~r.~-- zm.~1922~r.), Janiny\indeksosoby[Gomulska]{Janina}{Matuszewska}{1923}{Desno}
(ur.~1923~r.~-- zm.~1996~r.) oraz Marii\indeksosoby[Gomulska]{Maria}{Lewińska}{1925}{Desno}
(ur.~1925~r.~-- zm.~2013~r.), z~czego tylko trójka córek dożyła wieku dorosłego
-- potomkowie Wiktorii oraz Janiny mieszkają do~dnia dzisiejszego na~Deśnie.

\begin{figure}[!ht]
    \vspace*{0.5cm}
    \centering \includegraphics[width=0.87\linewidth]{
        1910_Marianna_Gomulska_Jan_Świtkowski_akt_ślubu_parafia_Długa_Kościelna_wpis_17.jpg}
    \captionsetup{format=hang}
    \caption{Akt ślubu Marianny Gomulskiej oraz Jana Świtkowskiego -- par. 
    Długa Kościelna 1910~rok (17/1910) \cite{par_dluga}.}
    \label{fig:mgomulska_1910}
\end{figure}

Pierwsza córka Stanisława i~Elżbiety Gomulskich,
\textbf{Marianna}\indeksosoby[Gomulska]{Marianna}{Świtkowska}{1884}{Jarząbek}[Sxyzwitkowska],
w~1910~roku wyszła za mąż za Jana Świtkowskiego\indeksosoby{Jan}{Świtkowski}{1885}{Cisie}[Sxyzwitkowski]
(ur. 1885~r. -- zm. 1912~r.) z~pobliskiej wsi Cisie (patrz ryc.
\ref{fig:mgomulska_1910}). Ich związek nie trwał długo, bo~już w~1912~roku Jan
zmarł, osieracając ich jedyne dziecko -- syna
Stanisława\indeksosoby{Stanisław}{Świtkowski}{1911}{Cisie}[Sxyzwitkowski]
(ur.~1911~r.~-- zm. 2011~r.). W~1918~roku Marianna wyszła ponownie za mąż,
co~ciekawe, za~młodszego brata Jana -- Władysława
Świtkowskiego\indeksosoby{Władysław}{Świtkowski}{1889}{Cisie}[Sxyzwitkowski]
(ur.~1889~r.~-- zm. 1951~r.) -- patrz ryc. \ref{fig:mgomulska_1918}. Para nie
doczekała się potomstwa.

\begin{figure}[!ht]
    \vspace*{0.5cm}
    \centering \includegraphics[width=0.9\linewidth]{
        1918_Marianna_Świtkowska_Gomulska_Władysław_Świtkowski_akt_ślubu_parafia_Długa_Kościelna_wpis_25.JPG}
    \captionsetup{format=hang}
    \caption{Akt ślubu Marianny Świtkowskiej oraz Władysława Świtkowskiego --
    par. Długa Kościelna 1918~rok (25/1918) \cite{par_dluga}.}
    \label{fig:mgomulska_1918}
\end{figure}

Druga córka Stanisława i~Elżbiety Gomulskich, \textbf{Antonina}, w~1909~roku
wyszła za~mąż za~Tomasza Forysia z~Aleksandrówki i~tam też osiedlili się
po~swoim ślubie (patrz ryc. \ref{fig:agomulska_1909}). Państwo Forysiowie
dochowali się trzech synów: Stanisława (ur. 1910~r. -- zm. 1988~r.), Czesława
(ur. 1914~r. -- zm. 1917~r.) oraz Tadeusza (ur. 1920~r. -- zm. 1994~r.).

\begin{figure}[!ht]
    \vspace*{0.5cm}
    \centering \includegraphics[width=0.87\linewidth]{
        1909_Antonina_Gomulska_Tomasz_Foryś_akt_ślubu_parafia_Długa_Kościelna_wpis_16.jpg}
    \captionsetup{format=hang}
    \caption{Akt ślubu Antoniny Gomulskiej oraz Tomasza Forysia --
    par. Długa Kościelna 1909~rok (16/1909) \cite{par_dluga}.}
    \label{fig:agomulska_1909}
\end{figure}
E
Drugi syn Stanisława i~Elżbiety Gomulskich,
\textbf{Józef}\indeksosoby{Józef}{Gomulski}{1892}{Anielina}, w~1918~roku
ożenił się ze swoją sąsiadką, Marcjanną Karłowską\indeksosoby[Karłowska]{Marcjanna}{Gomulska}{1899}{Desno}
(ur. 1920~r. -- zm. 1984~r.), córką
Andrzeja\indeksosoby{Andrzej}{Karłowski}{1877}{Cisie} (ur. 1877~r. --
zm.~1933~r.) i~Agnieszki\indeksosoby[Dębska]{Agnieszka}{Karłowska}{1881}{Zastawie}
z~domu Dębskiej (ur. 1881~r. -- zm. ?) Karłowskich, współudziałowców spółki
\emph{Desna}. Para po~ślubie zamieszkała na~Deśnie (patrz ryc.
\ref{fig:jgomulski_1918}). Józef i~Marcjanna Gomulscy dochowali się piątki
dzieci: Jana\indeksosoby{Jan}{Gomulski}{1918}{Desno} (ur. 1918~r. -- zm. ?),
Kazimiery\indeksosoby[Gomulska]{Kazimiera}{Łukawska}{1920}{Desno}[Lukawska],
(ur. 1920~r. -- zm. ?), Marianny\indeksosoby{Marianna}{Gomulska}{1924}{Desno}
(ur. 1924~r. -- zm. 1925~r.), Marianny\indeksosoby[Gomulska]{Marianna}{Kiełek}{1926}{Desno}
(ur. 1926~r. -- zm. 2022~r.), Romana\indeksosoby{Roman}{Gomulski}{1930}{Desno}
(ur. 1930~r. -- zm. 1930~r.) oraz Józefy\indeksosoby[Gomulska]{Józefa}{Pietrzak}{1932}{Desno}
(ur. 1932~r. -- zm.~?).

\begin{figure}[!ht]
    \vspace*{0.5cm}
    \centering \includegraphics[width=0.8\linewidth]{
        1918_Józef_Gomulski_Marcjanna_Karłowska_akt_ślubu_parafia_Długa_Kościelna_wpis_15.JPG}
    \captionsetup{format=hang}
    \caption{Akt ślubu Józefa Gomulskiego oraz Marcjanny Karłowskiej --
    par. Długa Kościelna 1918~rok (15/1918) \cite{par_dluga}.}
    \label{fig:jgomulski_1918}
\end{figure}

Ślub Józefa\indeksosoby{Józef}{Gomulski}{1892}{Anielina}
i~Marcjanny\indeksosoby[Gomulska]{Marianna}{Świtkowska}{1884}{Jarząbek}[Sxyzwitkowska]
Gomulskich był prawdopodobnie momentem, w~którym Stanisław i~Elżbieta Gomulscy
opuścili, wraz z~dwoma niezamężnymi jeszcze córkami, Stanisławą i~Katarzyną,
oraz najmłodszym synem Władysławem, wieś Desno, przeprowadzając się spowrotem
do~rodzinnej miejscowości Stanisława -- Anieliny. Przed przeprowadzką,
Stanisław zdążył przepisać na Józefa pozostałą mu część ziemi na~Deśnie oraz
udziały w~spółce \emph{Desna} (patrz ryc. \ref{fig:desna_uchwala}). Niedługo
po~przeprowadzce do~Anieliny, cała trójka dzieci Stanisława i~Elżbiety
Gomulskich wzięła śluby z~partnerami z~tamtych okolic, osiadając ostatecznie
razem z~rodzicami na~obrzeżach Mińska Mazowieckiego\footnote{Ich dalsza
historia została szerzej omówiona na~końcu podrozdziału poświęconego Anielinie.}.

\begin{figure}[!ht]
    \vspace*{0.4cm}
    \centering \includegraphics[width=0.78\linewidth]{
        1972_Jan_Gomulski_Eleonora_Gomulska_zdjęcie_nagrobka4_cmentarz_Długa_Kościelna.jpg}
    \captionsetup{format=hang}
    \caption{Zdjęcie grobu małżonków Jana i~Eleonory Gomulskich oraz członków
    ich najbliższej rodziny, znajdującego się na cmentarzu parafialnym parafii św. Anny
    w~Długiej Kościelnej (sektor: A3, rząd: 1, numer: 27)}
    \label{fig:jgomulska_1972}
\end{figure}

\begin{figure}[!ht]
    \vspace*{0.4cm}
    \centering \includegraphics[width=0.78\linewidth]{
        1951_Władysław_Świtkowski_Marianna_Świtkowska_zdjęcie_nagrobka1_cmentarz_Długa_Kościelna.jpg}
    \captionsetup{format=hang}
    \caption{Zdjęcie grobu małżonków Władysława oraz Marianny Świtkowskich
    znajdującego, się na cmentarzu parafialnym parafii św. Anny w~Długiej
    Kościelnej (sektor: B11, rząd: 11, numer: 11)}
    \label{fig:wswitkowski_1951}
\end{figure}

\begin{figure}[!ht]
    \vspace*{0.4cm}
    \centering \includegraphics[width=0.82\linewidth]{
        1985_Czesław_Mateuszewski_Janina_Matuszewska_zdjęcie_nagrobka_cmentarz_Długa_Kościelna.jpg}
    \captionsetup{format=hang}
    \caption{Zdjęcie grobu małżonków Czesława oraz Janiny Matuszewskich,
    znajdującego się na cmentarzu parafialnym parafii św. Anny w~Długiej
    Kościelnej (sektor: B3, rząd: 11, numer: 3)}
    \label{fig:cmatuszewski_1985}
\end{figure}

\begin{figure}[!ht]
    \vspace*{0.4cm}
    \centering \includegraphics[width=0.82\linewidth]{
        2013_Maria_Lewińska_zdjęcie_nagrobka_cmentarz_Długa_Kościelna.jpg}
    \captionsetup{format=hang}
    \caption{Zdjęcie grobu Marii Lewińskiej oraz jej córki, znajdującego się
    na cmentarzu parafialnym parafii św. Anny w~Długiej Kościelnej (sektor:
    A3, rząd: 1, numer: 28)}
    \label{fig:mlewinska_2012}
\end{figure}

\begin{figure}[!ht]
\vspace*{0.5cm}
\centering
\begin{tikzpicture}[
  % Definiujemy stałe odstępy
  node distance=1cm and 1cm, 
  every node/.style={draw, rectangle, align=center, minimum width=3.8cm,
  inner sep=5pt}]

% --- POKOLENIE 0 ---
\node (Osiowa0) [draw=none, minimum width=0pt, inner sep=0pt] at (0,0) {};
\node (Antoni) [left=1.5cm of Osiowa0] {Antoni Gomoła \\ {
    \footnotesize \emph{Wnętrzne, 1764 r.}}};
\node (Petronella) [right=1.5cm of Osiowa0] {Petronella Sokół \\ {
    \footnotesize \emph{Mysłów, 1768 r.}}};

\draw[-] (Antoni) -- (Petronella)
    node[midway, below, draw=none, minimum width=0pt,
    font=\scriptsize] (Slub0) {\emph{Wilczyska, 1783 r.}};

% --- POKOLENIE 1 ---
\node (Osiowa1) [draw=none, minimum width=0pt, inner sep=0pt,
below=1.5cm of Slub0] {};
\node (Jan) [left=1.5cm of Osiowa1] {Jan Gomoła \\ {
    \footnotesize \emph{Mysłów, 1789 r.}}};
\node (Łucja) [right=1.5cm of Osiowa1] {Łucja Rokicka \\ {
    \footnotesize \emph{Rudzienko, 1790 r.}}};

\draw[-] (Jan) -- (Łucja)
    node[midway, above, draw=none, minimum width=0pt,
    font=\scriptsize] (Slub1_Gora) { }
    node[midway, below, draw=none, minimum width=0pt,
    font=\scriptsize] (Slub1) {\emph{Kołbiel, 1814 r.}};

\draw[->] (Slub0) -- (Slub1_Gora);

% --- POKOLENIE 2 ---
\node (Osiowa2) [draw=none, minimum width=0pt, inner sep=0pt,
below=1.5cm of Slub1] {};
\node (Piotr) [left=1.5cm of Osiowa2] {Piotr Franciszek Gumulski \\ {
    \footnotesize \emph{Rudzienko, 1815 r.}}};
\node (Marianna) [right=1.5cm of Osiowa2] {Marianna Chłopik \\ {
    \footnotesize \emph{Zamienie, 1819 r.}}};
    
\draw[-] (Piotr) -- (Marianna)
    node[midway, above, draw=none, minimum width=0pt,
    font=\scriptsize] (Slub2_Gora) { }
    node[midway, below, draw=none, minimum width=0pt,
    font=\scriptsize] (Slub2) {\emph{Mińsk, 1835 r.}};

\draw[->] (Slub1) -- (Slub2_Gora);

% --- POKOLENIE 3 ---
\node (Osiowa3) [draw=none, minimum width=0pt, inner sep=0pt,
below=1.5cm of  Slub2] {};
\node (Stanislaw) [left=1.5cm of Osiowa3] {Stanisław Gomulski \\ {
    \footnotesize \emph{Anielina, 1863 r.}}};
\node (Elzbieta) [right=1.5cm of Osiowa3] {Elżbieta Brodzińska \\ {
    \footnotesize \emph{Grębiszew, 1857 r.}}};

\draw[-] (Stanislaw) -- (Elzbieta)
    node[midway, above, draw=none, minimum width=0pt,
    font=\scriptsize] (Slub3_Gora) { }
    node[midway, below, draw=none, minimum width=0pt,
    font=\scriptsize] (Slub3) {\emph{Mińsk, 1882 r.}};

\draw[->] (Slub2) -- (Slub3_Gora);

\end{tikzpicture}
\end{figure}

Stanisław Gomulski\indeksosoby{Stanisław}{Gomulski}{1863}{Anielina} zmarł
15~grudnia 1929~roku, w~wieku 66 lat, w~Anielinie,
jego żona Elżbieta\indeksosoby[Brodzińska]{Elżbieta}{Gomulska}{1857}{Grębiszew}
zmarła półtora roku przed nim -- 12 lipca 1928~roku. Chociaż Stanisław
i~Elżbieta Gomulscy zmarli po około 10~latach od~powrotu do~Anieliny
i~prawdopodobnie zostali pochowani na~cmentarzu parafialnym parafii
pod wezwaniem Narodzenia Najświętszej Maryi Panny w~Mińsku Mazowieckim, to
istotna część ich rodziny dożyła swoich dni na Deśnie lub w~okolicznych
miejscowościach -- ich groby można odnaleźć na cmentarzu parafialnym
parafii św. Anny w~Długiej Kościelnej. Ryciny: \ref{fig:jgomulska_1972},
\ref{fig:wswitkowski_1951}, \ref{fig:cmatuszewski_1985} oraz
\ref{fig:mlewinska_2012} przedstawiają przykłady niektórych z~ich nagrobków.
Bezpośredni potomkowie Stanisława i~Elżbiety Gomulskich, mieszkają na terenie
obecnego Desna do dnia dzisiejszego -- żaden z~nich nie nosi już wprawdzie
nazwiska Gomulski, ze względu na silną dominację kobiet w~tej gałęzi rodziny,
ale noszą nazwiska Zawada oraz Matuszewski.

\newpage
\ifodd\value{page}\hbox{}\newpage\fi

\section*{Potomkowie Stanisława i~Julianny Gomulskich}

Ostatnim członkiem rodziny Gomulskich, który osiedlił się na terenie obecnego
Desna pod koniec XIX~wieku, był
Stanisław\indeksosoby{Stanisław}{Gomulski}{1861}{Wólka Mińska} Gomulski
(ur. 1861~r. -- zm. 1910~r.), drugi syn Jana Gomulskiego
i~Agnieszki\indeksosoby[Piotrkowicz]{Agnieszka}{Gomulska}{1840}{Budy Barcząckie}
z~domu Piotrkowicz, wnuk Piotra Franciszka Gomulskiego i~Marianny z~domu
Chłopik. Stanisław Gomulski 19~listopada 1881~roku wziął ślub w~parafii pod
wezwaniem Narodzenia Najświętszej Maryi Panny w~Mińsku Mazowieckim z~Julianną
Gągol\indeksosoby[Gągol]{Julianna}{Gomulska}{1862}{Dłużka} (ur. 1862~r. -- zm. ?) --
patrz ryc. \ref{fig:sgomulski_1881}. Julianna urodziła się 12~kwietnia
1862~roku we wsi Dłużka (pod Mińskiem Mazowieckim), była córką
Józefa\indeksosoby{Józef}{Gągol}{1835}{Dłużka} Gągola (ur.~1835~r.~--
zm.~1866~r.) oraz Franciszki\indeksosoby[Kaszewska]{Franciszka}{Gągol}{1837}{Karolina}
z~domu Kaszewskiej (ur.~1835~r.~-- zm.~1860~r.) -- patrz ryc.
\ref{fig:jgagol_1862}. Państwo Gomulscy przez pierwsze 16~lat po ślubie
mieszkali w~Wólce Mińskiej, rodzinnej wsi Stanisława, gdzie urodziła im~się
szóstka najstarszych dzieci:
Piotr\indeksosoby{Piotr}{Gomulski}{1884}{Wólka Mińska},
Franciszka\indeksosoby[Gomulska]{Franciszka}{Jackiewicz}{1886}{Wólka Mińska},
Marianna\indeksosoby[Gomulska]{Marianna}{Szczepańska}{1889}{Wólka Mińska},
Karolina\indeksosoby[Gomulska][Piotrkowicz]{Karolina}{Siporska}{1891}{Wólka Mińska},
Bezimienny\indeksosoby{Bezimienny}{Gomulski}{1893}{Wólka Mińska} oraz
Józef\indeksosoby{Józef}{Gomulski}{1895}{Wólka Mińska}.

\begin{figure}[!ht]
    \vspace*{0.5cm}
    \centering \includegraphics[width=0.9\linewidth]{
        1862_Julianna_Gomulska_Gągol_akt_chrztu_parafia_Mińsk_Mazowiecki_wpis_118.jpg}
    \captionsetup{format=hang}
    \caption{Akt chrztu Julianny Gągol -- par. Mińsk Mazowiecki 1862~rok (118/1862) 
    \cite{par_minsk2}.}
    \label{fig:jgagol_1862}
\end{figure}

Stanisław\indeksosoby{Stanisław}{Gomulski}{1861}{Wólka Mińska}
i~Julianna\indeksosoby[Gągol]{Julianna}{Gomulska}{1862}{Dłużka} Gomulscy
dochowali się łącznie dziesiątki dzieci, z~czego aż~dziewiątka dożyła wieku
dorosłego\footnote{Drzewo genealogiczne najbliższej rodziny Stanisława
i~Julianny Gomulskich zostało zaprezentowane
w~\hyperref[sec:drzewo_sgomulski2]{załączniku numer I} do~niniejszej książki.}:

\begin{itemize}
    \item Piotr Gomulski\indeksosoby{Piotr}{Gomulski}{1884}{Wólka Mińska}
    (ur. 1884~r. -- zm. 1910~r.) -- patrz: ryc. \ref{fig:pgomulski_1884},
    \item Franciszka Gomulska\indeksosoby[Gomulska]{Franciszka}{Jackiewicz}{1886}{Wólka Mińska}
    (ur. 1886~r. -- zm. ?) -- patrz: ryc. \ref{fig:fgomulska_1886},
    \item Marianna Gomulska\indeksosoby[Gomulska]{Marianna}{Szczepańska}{1889}{Wólka Mińska}
    (ur. 1889~r. -- zm. 1910~r.) -- patrz: ryc. \ref{fig:mgomulska_1889},
    \item Karolina Gomulska\indeksosoby[Gomulska][Piotrkowicz]{Karolina}{Siporska}{1891}{Wólka Mińska}
    (ur. 1891~r. -- zm. 1959~r.) -- patrz: ryc. \ref{fig:kgomulska_1891},
    \item Bezimienny Gomulski\indeksosoby{Bezimienny}{Gomulski}{1893}{Wólka Mińska}
    (ur. 1893~r. -- zm. 1893~r.) -- patrz: ryc. \ref{fig:bgomulski_1893},
    \item Józef Gomulski\indeksosoby{Józef}{Gomulski}{1895}{Wólka Mińska}
    (ur. 1895~r. -- zm. 1950~r.) -- patrz: ryc. \ref{fig:jgomulski_1895},
    \item Franciszek Gomulski\indeksosoby{Franciszek}{Gomulski}{1897}{Desno}
    (ur. 1897~r. -- zm. ?) -- patrz: ryc. \ref{fig:fgomulski_1897},
    \item Tomasz Gomulski\indeksosoby{Tomasz}{Gomulski}{1900}{Desno}
    (ur. 1900~r. -- zm. 1982~r.) -- patrz: ryc. \ref{fig:tgomulski_1900},
    \item Wincenty Gomulski\indeksosoby{Wincenty}{Gomulski}{1903}{Desno}
    (ur. 1903~r. -- zm. 1983~r.) -- patrz: ryc. \ref{fig:wgomulski_1903},
    \item Aleksander Gomulski\indeksosoby{Aleksander}{Gomulski}{1906}{Desno}
    (ur. 1906~r. -- zm. 1987~r.) -- patrz: ryc. \ref{fig:agomulski_1906}.
\end{itemize}

\begin{figure}[!ht]
    \vspace*{0.5cm}
    \centering \includegraphics[width=0.85\linewidth]{
        1881_Stanisław_Gomulski_Julianna_Gągol_akt_ślubu_parafia_Mińsk_Mazowiecki_wpis_56.jpg}
    \captionsetup{format=hang}
    \caption{Akt ślubu Stanisława Gomulskiego oraz Julianny Gągol --
    par. Mińsk Mazowiecki 1881~rok (56/1881) \cite{par_minsk2}.}
    \label{fig:sgomulski_1881}
\end{figure}

\begin{figure}[!ht]
    \vspace*{0.5cm}
    \centering \includegraphics[width=0.75\linewidth]{
        1884_Piotr_Gomulski_akt_chrztu_parafia_Mińsk_Mazowiecki_wpis_101.jpg}
    \captionsetup{format=hang}
    \caption{Akt chrztu Piotra Gomulskiego -- par. Mińsk Mazowiecki 1884~rok
    (101/1884)  \cite{par_minsk2}.}
    \label{fig:pgomulski_1884}
\end{figure}

\begin{figure}[!ht]
    \vspace*{0.5cm}
    \centering \includegraphics[width=0.75\linewidth]{
        1886_Franciszka_Jackowicz_Gomulska_akt_chrztu_parafia_Mińsk_Mazowiecki_wpis_120.jpg}
    \captionsetup{format=hang}
    \caption{Akt chrztu Franciszki Gomulskiej -- par. Mińsk Mazowiecki
    1886~rok (120/1886)  \cite{par_minsk2}.}
    \label{fig:fgomulska_1886}
\end{figure}

\begin{figure}[!ht]
    \vspace*{0.5cm}
    \centering \includegraphics[width=0.75\linewidth]{
        1889_Marianna_Szczepańska_Gomulska_akt_chrztu_parafia_Mińsk_Mazowiecki_wpis_9.jpg}
    \captionsetup{format=hang}
    \caption{Akt chrztu Marianny Gomulskiej -- par. Mińsk Mazowiecki
    1889~rok (9/1889)  \cite{par_minsk2}.}
    \label{fig:mgomulska_1889}
\end{figure}

\begin{figure}[!ht]
    \vspace*{0.5cm}
    \centering \includegraphics[width=0.75\linewidth]{
        1891_Karolina_Siporska_Piotrkowicz_Gomulska_akt_chrztu_parafia_Mińsk_Mazowiecki_wpis_444.jpg}
    \captionsetup{format=hang}
    \caption{Akt chrztu Karoliny Gomulskiej -- par. Mińsk Mazowiecki
    1891~rok (444/1891)  \cite{par_minsk2}.}
    \label{fig:kgomulska_1891}
\end{figure}

\begin{figure}[!ht]
    \vspace*{0.5cm}
    \centering \includegraphics[width=0.8\linewidth]{
        1893_Bezimienny_Gomulski_akt_chrztu_akt_śmierci_parafia_Mińsk_Mazowiecki_wpis_29.jpg}
    \captionsetup{format=hang}
    \caption{Akt chrztu Bezimiennego Gomulskiego -- par. Mińsk Mazowiecki
    1893~rok (29/1893)  \cite{par_minsk2}.}
    \label{fig:bgomulski_1893}
\end{figure}

\begin{figure}[!ht]
    \vspace*{0.5cm}
    \centering \includegraphics[width=0.83\linewidth]{
        1895_Józef_Gomulski_akt_chrztu_parafia_Mińsk_Mazowiecki_wpis_109.jpg}
    \captionsetup{format=hang}
    \caption{Akt chrztu Józefa Gomulskiego -- par. Mińsk Mazowiecki
    1895~rok (109/1895)  \cite{par_minsk2}.}
    \label{fig:jgomulski_1895}
\end{figure}

\begin{figure}[!ht]
    \vspace*{0.5cm}
    \centering \includegraphics[width=0.7\linewidth]{
        1897_Franciszek_Gomulski_akt_chrztu_parafia_Długa_Kościelna_wpis_102.jpg}
    \captionsetup{format=hang}
    \caption{Akt chrztu Franciszka Gomulskiego -- par. Mińsk Mazowiecki
    1897~rok (109/1897)  \cite{par_dluga}.}
    \label{fig:fgomulski_1897}
\end{figure}

\begin{figure}[!ht]
    \vspace*{0.5cm}
    \centering \includegraphics[width=0.7\linewidth]{
        1900_Tomasz_Gomulski_akt_chrztu_parafia_Długa_Kościelna_wpis_125.jpg}
    \captionsetup{format=hang}
    \caption{Akt chrztu Tomasza Gomulskiego -- par. Długa Kościelna
    1900~rok (125/1900)  \cite{par_dluga}.}
    \label{fig:tgomulski_1900}
\end{figure}

\begin{figure}[!ht]
    \vspace*{0.5cm}
    \centering \includegraphics[width=0.8\linewidth]{
        1903_Wincenty_Gomulski_akt_chrztu_parafia_Długa_Kościelna_wpis_11.jpg}
    \captionsetup{format=hang}
    \caption{Akt chrztu Wincentego Gomulskiego -- par. Długa Kościelna
    1903~rok (11/1903)  \cite{par_dluga}.}
    \label{fig:wgomulski_1903}
\end{figure}

\begin{figure}[!ht]
    \vspace*{0.5cm}
    \centering \includegraphics[width=0.8\linewidth]{
        1906_Aleksander_Gomulski_akt_chrztu_parafia_Długa_Kościelna_wpis_17.jpg}
    \captionsetup{format=hang}
    \caption{Akt chrztu Aleksandra Gomulskiego -- par. Długa Kościelna
    1906~rok (17/1906)  \cite{par_dluga}.}
    \label{fig:agomulski_1906}
\end{figure}

Około 1897~roku Stanisław\indeksosoby{Stanisław}{Gomulski}{1861}{Wólka Mińska}
i~Julianna\indeksosoby[Gągol]{Julianna}{Gomulska}{1862}{Dłużka} Gomulscy, wraz
z~szóstką dzieci, przeprowadzili się do oddalonej o~około 11~kilometrów
na~północny zachód od~Wólki Mińskiej wsi \emph{Desna}, gdzie przyszła na~świat
czwórka ich najmłodszych synów:
Franciszek\indeksosoby{Franciszek}{Gomulski}{1897}{Desno},
Tomasz\indeksosoby{Tomasz}{Gomulski}{1900}{Desno},
Wincenty\indeksosoby{Wincenty}{Gomulski}{1903}{Desno} oraz
Aleksander\indeksosoby{Aleksander}{Gomulski}{1906}{Desno}.
O~tym, że~małżonkowie Gomulscy sprowadzili się na~tereny obecnego Desna
dopiero około 1897~roku świadczy fakt, iż~Stanisław nie był wymieniony
na~liście pierwotnych współwłaścicieli spółki \emph{Desna} z~1895~roku
(w~przeciwieństwie do~jego stryjów Michała i~Stanisława), a~według akt parafii
św.~Anny w~Długiej Kościelnej w~październiku 1897~roku zgłosił on~urodzenie
syna Franciszka jako gospodarz zamieszkujący we wsi \emph{Desna} (patrz: ryc.
\ref{fig:fgomulski_1897}). W~księgach hipotecznych dóbr \emph{Desna} widnieje
natomiast informacja, iż Stanisław Gomulski 6~października 1900~roku odkupił
udziały w~spółce \emph{Desna} za kwotę 600~rubli od~Konstatntego
Kaima\indeksosoby{Konstanty}{Kaim}{1860}{Stanisławów} przejmując w~ten sposób
własność 8~hektarów ziemi rolnej oraz dług na~kwotę około 500~rubli (patrz:
ryc. \ref{fig:sgomulski_1901}).

\begin{figure}[!ht]
    \vspace*{0.5cm}
    \centering \includegraphics[width=0.62\linewidth]{
        Stanislaw_Gomulski_zakup_dzialki_Desna.JPG}
    \captionsetup{format=hang}
    \caption{Wpis w~księgach hipotecznych dóbr \emph{Desna} świadczący
    o~zakupie przez Stanisława Gomulskiego 8~hektarów ziemi rolnej na terenie
    wsi \emph{Desna} \cite{desna1}.}
    \label{fig:sgomulski_1901}
\end{figure}

Spośród dziewiątki dzieci małżonków Gomulskich, tylko czwórka synów osiedliła
się na~stałe na terenie obecnego Desna:
Józef\indeksosoby{Józef}{Gomulski}{1895}{Wólka Mińska},
Tomasz\indeksosoby{Tomasz}{Gomulski}{1900}{Desno},
Wincenty\indeksosoby{Wincenty}{Gomulski}{1903}{Desno} oraz
Aleksander\indeksosoby{Aleksander}{Gomulski}{1906}{Desno} -- ich potomkowie
żyją na~Deśnie do~dnia dzisiejszego\footnote{Jednym z~nich jest autor
niniejszej ksiażki.}.

\begin{figure}[!ht]
    \vspace*{0.5cm}
    \centering \includegraphics[width=0.8\linewidth]{
        1910_Piotr_Gomulski_Marcjanna_Ostrowska_akt_ślubu_parafia_Długa_Kościelna_wpis_8.jpg}
    \captionsetup{format=hang}
    \caption{Akt ślubu Piotra Gomulskiego oraz Marcjanny Ostrowskiej --
    par. Długa Kościelna 1910~rok (8/1910) \cite{par_dluga}.}
    \label{fig:pgomulski_1910}
\end{figure}

Najstarszy syn Stanisława i~Julianny Gomulskich,
\textbf{Piotr}\indeksosoby{Piotr}{Gomulski}{1884}{Wólka Mińska}, 24~stycznia
1910~roku, w~Długiej Kościelnej, ożenił się z~Marcjanną
Ostrowską\indeksosoby[Ostrowska][Gomulska]{Marcjanna}{Moskwik}{1887}{Mrowiska}
(ur. 1887~r. -- zm. 1940~r.), panną pochodzącą z~pobliskich Mrowisk. Ich
związek nie trwał długo, bo~nieco ponad dwa miesiące po~ich ślubie -- pod
koniec marca 1910~roku, Piotr Gomulski zmarł. Para nie dochowała się potomstwa,
a~Marcjanna dwa lata po~śmierci swojego pierwszego męża wyszła ponownie za~mąż
za~Franciszka Moskwika\indeksosoby{Franciszek}{Moskwik}{1889}{Żurawka}
(ur. 1889~r. -- zm. 1973~r.) z~Żurawki.

\begin{figure}[!ht]
    \vspace*{0.5cm}
    \centering \includegraphics[width=0.75\linewidth]{
        1905_Franciszka_Gomulska_Franciszek_Jackiewicz_akt_ślubu_parafia_Długa_Kościelna_wpis_12.jpg}
    \captionsetup{format=hang}
    \caption{Akt ślubu Franciszki Gomulskiej oraz Franciszka Jackiewicza --
    par. Długa Kościelna 1905~rok (12/1905) \cite{par_dluga}.}
    \label{fig:fgomulska_1905}
\end{figure}

Najstarsza córka Stanisława i~Julianny Gomulskich,
\textbf{Franciszka}\indeksosoby[Gomulska]{Franciszka}{Jackiewicz}{1886}{Wólka Mińska},
w~marcu 1905~roku, w~Długiej Kościelnej, wyszła za mąż za Franciszka
Jackiewicza\indeksosoby{Franciszek}{Jackiewicz}{1874}{Poręby}
(ur. 1874~r. -- zm. ?), kawalera pochodzącego z~Poręb. Po ślubie para
zamieszkała w~rodzinnej wsi Franciszka, gdzie dochowała się szóstki dzieci:
Juliana\indeksosoby{Julian}{Jackiewicz}{1906}{Poręby} (ur. 1906~r. -- zm. 1979~r.),
Józefa\indeksosoby{Józef}{Jackiewicz}{1909}{Poręby} (ur. 1909~r. --zm. 1990~r.),
Jana\indeksosoby{Jan}{Jackiewicz}{?}{Poręby} (ur. ? -- zm. ?),
Marianny\indeksosoby{Marianna}{Jackiewicz}{?}{Poręby} (ur. ? -- zm. ?),
Mariana\indeksosoby{Marian}{Jackiewicz}{?}{Poręby} (ur. ? -- zm. ?) oraz
Edwarda\indeksosoby{Edward}{Jackiewicz}{?}{Poręby} (ur. ? -- zm. ?).

\begin{figure}[!ht]
    \vspace*{0.5cm}
    \centering \includegraphics[width=0.77\linewidth]{
        1908_Marianna_Gomulska_Jan_Szczepański_akt_ślubu_parafia_Długa_Kościelna_wpis_12.jpg}
    \captionsetup{format=hang}
    \caption{Akt ślubu Marianny Gomulskiej oraz Jana Szczepańskiego --
    par. Długa Kościelna 1908~rok (12/1908) \cite{par_dluga}.}
    \label{fig:mgomulska_1908}
\end{figure}

Druga córka Stanisława i~Julianny Gomulskich, 
\textbf{Marianna}\indeksosoby[Gomulska]{Marianna}{Szczepańska}{1889}{Wólka Mińska},
w~lutym 1908~roku, w~Długiej Kościelnej, wyszła za~mąż za~Jana
Szczepańskiego\indeksosoby{Jan}{Szczepański}{1879}{Pustelnik} (ur.
1879~r. -- zm. 1941~r.), kawalera pochodzącego z~Pustelnika. Państwo
Szczepańscy nie dochowali się potomstwa, gdyż niewiele ponad dwa lata po~ich
ślubie, w~kwietniu 1910~roku Marianna zmarła w~rodzinnym Deśnie, w~wieku
zaledwie 21~lat.

\begin{figure}[!ht]
    \vspace*{0.5cm}
    \centering \includegraphics[width=0.82\linewidth]{
        1910_Piotr_Gomulski_akt_zgonu_parafia_Długa_Kościelna_wpis_18.jpg}
    \captionsetup{format=hang}
    \caption{Akt zgonu Piotra Gomulskiego -- par. Długa Kościelna 1908~rok
    (18/1908) \cite{par_dluga}.}
    \label{fig:pgomulski2_1910}
\end{figure}

\begin{figure}[!ht]
    \vspace*{0.5cm}
    \centering \includegraphics[width=0.82\linewidth]{
        1910_Stanisław_Gomulski_akt_zgonu_parafia_Długa_Kościelna_wpis_19.jpg}
    \captionsetup{format=hang}
    \caption{Akt zgonu Stanisława Gomulskiego -- par. Długa Kościelna 1910~rok
    (19/1910) \cite{par_dluga}.}
    \label{fig:sgomulski_1910}
\end{figure}

\begin{figure}[!ht]
    \vspace*{0.5cm}
    \centering \includegraphics[width=0.93\linewidth]{
        1910_Marianna_Szczepańska_Gomulska_akt_zgonu_parafia_Długa_Kościelna_wpis_20.jpg}
    \captionsetup{format=hang}
    \caption{Akt zgonu Marianny Szczepańskiej z~domu Gomulskiej -- par. Długa
    Kościelna 1910~rok (20/1910) \cite{par_dluga}.}
    \label{fig:mgomulska2_1910}
\end{figure}

Rok 1910~można bez wahania nazwać prawdziwym \emph{annus horribilis} dla
rodziny Stanisława\indeksosoby{Stanisław}{Gomulski}{1861}{Wólka Mińska}
i~Julianny\indeksosoby[Gągol]{Julianna}{Gomulska}{1862}{Dłużka} Gomulskich.
W~ciągu zaledwie ośmiu dni, następujących bezpośrednio
po~Wielkanocy\footnote{Wielka Niedziela w~1910~roku wypadała 27~marca.},
zmarły aż~trzy dorosłe osoby z~tej rodziny — niemal jedna piąta najbliższych
Stanisława i~Julianny:

\begin{itemize}
\item W~Wielki Poniedziałek, 28~marca 1910~roku, osiem dni po swoich
26.~urodzinach, zmarł Piotr Gomulski\indeksosoby{Piotr}{Gomulski}{1884}{Wólka Mińska},
najstarszy syn Stanisława i~Julianny (patrz: ryc. \ref{fig:pgomulski2_1910}).
Piotr pozostawił po~sobie żonę Marcjannę, z~którą wziął ślub nieco ponad dwa
miesiące wcześniej.
\item Tydzień po Wielkanocy, 3~kwietnia~1910~roku, zmarł w~wieku 49~lat,
sam Stanisław Gomulski\indeksosoby{Stanisław}{Gomulski}{1861}{Wólka Mińska}
(patrz: ryc. \ref{fig:sgomulski_1910}), ojciec Piotra i~Marianny. Osierocił
on~piątkę nieletnich synów, z~których najmłodszy, Aleksander, miał zaledwie
cztery lata.
\item Osiem dni po Wielkanocy, 4~kwietnia 1910~roku, w~wieku zaledwie
21~lat, zmarła Marianna
Szczepańska\indeksosoby[Gomulska]{Marianna}{Szczepańska}{1889}{Wólka Mińska}
z~domu Gomulska, młodsza siostra Piotra i~córka Stanisława (patrz: ryc.
\ref{fig:mgomulska2_1910}). Marianna pozostawiła po~sobie męża Jana, z~ktorym
wzięła ślub zaledwie dwa lata wcześniej.
\end{itemize}

W~ich aktach zgonu nie sposób jest odnaleźć informacji dotyczącej przyczyny
śmierci żadnego z~nich, być może to~jakaś zaraźliwa choroba (zapalenie płuc
lub dur brzuszny) przywieziona przez kogoś na Wielkanoc uśmierciła ludzi w~tak
młodym wieku. Jedne co wiadomo to fakt, iż wszyscy troje zmarli na Deśnie,
były to trzy kolejne zgony odnotowane w~księgach metrykalnych parafii św. Anny
w~Długiej Kościelnej w~roku 1910 -- pod numerami 18, 19 oraz 20. Wszyscy troje
zostali pochowali we wspólnym grobie, który istnieje do~dnia dzisiejszego na
cmentarzu parafialnym w~Długiej Kościelnej\footnote{Być może jest to jeden
z~najstarszych grobów na tym cmentarzu, co do którego można mieć pewność kiedy
dokładnie został utworzony.}, choć czas wytarł już część napisów na nagrobku
-- patrz: ryc. \ref{fig:grob_gomulskich_1910}.

\begin{figure}[!ht]
    \vspace*{0.4cm}
    \centering \includegraphics[width=1.0\linewidth]{
        1910_Stanisław_Gomulski_Julianna_Gomulska_zdjęcie_nagrobka1_cmentarz_Długa_Kościelna.jpg}
    \captionsetup{format=hang}
    \caption{Zdjęcie grobu najbliższej rodziny Stanisława i~Julianny
    Gomulskich, znajdującego się od~1910~roku na cmentarzu parafialnym parafii
    św. Anny w~Długiej Kościelnej (sektor: A2, rząd: 8, numer: 13)}
    \label{fig:grob_gomulskich_1910}
\end{figure}

W~1915~roku, w~pięć lat po śmierci swojego ojca, brata i~siostry,
dwudziestoczteroletnia wówczas
\textbf{Karolina Gomulska}\indeksosoby[Gomulska][Piotrkowicz]{Karolina}{Siporska}{1891}{Wólka Mińska},
wzięła ślub z~68-letnim Janem
Piotrkowiczem\indeksosoby{Jan}{Piotrkowicz}{1847}{Kąty Goździejewskie}
(ur. 1847~r. -- zm.~?), największym niegdyś udziałowcem spółki
\emph{Desna}\footnote{Do~momentu, gdy około 1904~roku, nie przepisał
on~większości posiadanej ziemi w~\emph{Deśnie} swojemu synowi Andrzejowi
i~jego żonie Magdalenie.} (patrz: ryc. \ref{fig:kgomulska_1915}).
Jan Piotrkowicz był już wówczas wdowcem po swojej żonie Krystynie (z~domu
Link), która zmarła około pół roku wcześniej.

\begin{figure}[!ht]
    \vspace*{0.5cm}
    \centering \includegraphics[width=0.98\linewidth]{
        1915_Karolina_Gomulska_Jan_Piotrkowicz_akt_ślubu_parafia_Długa_Kościelna_wpis_12.jpg}
    \captionsetup{format=hang}
    \caption{Akt ślubu Karoliny Gomulskiej oraz Jana Piotrkowicza -- par.
    Długa Kościelna 1915~rok (12/1915) \cite{par_dluga}.}
    \label{fig:kgomulska_1915}
\end{figure}

Świeżoupieczonych małżonków Piotrkowiczów dzieliły w~momencie ślubu aż 44~lata
różnicy wieku, stąd też prawdopodobnie nie dochowali się oni potomstwa. Jan
Piotrkowicz zmarł po~około 10-latach nowego małżeństwa, o~czym świadczy wpis
w~księdze hipotecznej wsi \emph{Desna} z~września 1926~roku informujący
o~zakupie 3/10~udziału w~działce numer~1a~-- Karolina z~Gomulskich Piotrkowicz
została określona w~tym wpisie jako wdową mieszkającą we wsi Olesin --
patrz: ryc. \ref{fig:gomulscy_zakup}\footnote{Co ciekawe, pod przytoczonym
wpisem widnieje własnoręczny podpis Karoliny
Piotrkowicz\indeksosoby[Gomulska][Piotrkowicz]{Karolina}{Siporska}{1891}{Wólka Mińska},
co~stoi w~sprzecznosci z~rodzinną legendą autora niniejszej książki, która
głosiła, iż~ciotka Karólka miała problemy z~pisaniem i~żeby udowodnić swoim
trzem najmłodszym braciom:
Tomkowi\indeksosoby{Tomasz}{Gomulski}{1900}{Desno},
Wickowi\indeksosoby{Wincenty}{Gomulski}{1903}{Desno}
i~Olkowi\indeksosoby{Aleksander}{Gomulski}{1906}{Desno}, że~portafi pisać,
podpisała się nieumyślnie pod aktem przekazania im~całego swojego majątku.}.

\begin{figure}[!ht]
    \centering
    \includegraphics[width=0.48\linewidth]{Gomulscy_zakup_dzialki_1.JPG}
    \hfill
    \includegraphics[width=0.48\linewidth]{Gomulscy_zakup_dzialki_3.JPG}
    \captionsetup{format=hang}
    \caption{Pierwsza i~ostatnia strona wpisu w~księdze hipotecznej wsi
    \emph{Desna} dotyczącego zakupu 3/10~udziału w~działce numer 1a
    \cite{desna1}.}
    \label{fig:gomulscy_zakup}
\end{figure}

Kilkanaście lat po śmierci Jana Piotrkowicza\indeksosoby{Jan}{Piotrkowicz}{1847}{Kąty Goździejewskie},
Karolina Piotrkowicz\indeksosoby[Gomulska][Piotrkowicz]{Karolina}{Siporska}{1891}{Wólka Mińska},
w~lipcu 1943~roku w~kościele parafii św. Anny w~Długiej Kościelnej, wyszła ponownie
za~mąż -- za~kolejnego wdowca, tym razem za~62-letniego Jana
Siporskiego\indeksosoby{Jan}{Siporski}{1881}{Chobot} (ur.~1881~r.~-- zm. ?)
z~pobliskiej wsi Chobot (patrz: ryc.\ref{fig:kgomulska_1943}). Karolina miała
wówczas 52~lata, była więc młodsza o~jedynie 10~lat od~swojego nowego męża --
para ze~względu na~zaawansowany wiek nie doczekała się dzieci, ale przeżyła
wspólnie około 16~lat, do momentu śmierci Karoliny w~1959 roku. Karolina
Siporska z~domu Gomulska została pochowana w grobie rodzinnym swoich rodziców
-- patrz: ryc. \ref{fig:grob_gomulskich_1910}.

\begin{figure}[!ht]
    \vspace*{0.5cm}
    \centering \includegraphics[width=0.9\linewidth]{
        1943_Karolina_Piotrkowicz_Gomulska_Jan_Siporski_akt_ślubu_parafia_Długa_Kościelna_wpis_17.jpg}
    \captionsetup{format=hang}
    \caption{Akt ślubu Karoliny Piotrkowicz z~domu Gomulskiej oraz Jana
    Siporskiego -- par. Długa Kościelna 1943~rok (17/1943) \cite{par_dluga}.}
    \label{fig:kgomulska_1943}
\end{figure}

Trzeci syn Stanisława i~Julianny Gomulskich,
\textbf{Józef}\indeksosoby{Józef}{Gomulski}{1895}{Wólka Mińska}, w~październiku
1916~roku, w~Długiej Kościelnej, wziął ślub z~Feliksą
Tkaczyk\indeksosoby[Tkaczyk]{Feliksa}{Gomulska}{1897}{Desno}\footnote{Józef
i~Feliksa Gomulscy to pradziadkowie autor niniejszej książki -- drzewo
genealogiczne ich najbliższej rodziny zostało zaprezentowane
w~\hyperref[sec:drzewo_jzfgomulski]{załączniku numer I}~do niniejszej książki.}
(ur. 1897~r. -- zm. 1974~r.), panną pochodzącą z~Desna, córką
Franciszka\indeksosoby{Franciszek}{Tkaczyk}{1864}{Olesin} (ur.~1864~r.~--
zm.~1912~r.) oraz Teofili\indeksosoby[Konowrocka]{Teofila}{Tkaczyk}{1862}{Chobot}
z~domu Konowrockiej Tkaczków (ur.~1864~r.~-- zm~1912~r.) współzałożycieli
spółki \emph{Desna} (patrz: ryc. \ref{fig:jgomulski_1916}).

\begin{figure}[!ht]
    \vspace*{0.5cm}
    \centering \includegraphics[width=0.82\linewidth]{
        1916_Józef_Gomulski_Feliksa_Tkaczyk_akt_ślubu_parafia_Długa_Kościelna_wpis_21.jpg}
    \captionsetup{format=hang}
    \caption{Akt ślubu Józefa Gomulskiego oraz Feliksy Tkaczyk -- par. Długa
    Kościelna 1916~rok (21/1916) \cite{par_dluga}.}
    \label{fig:jgomulski_1916}
\end{figure}

Józef\indeksosoby{Józef}{Gomulski}{1895}{Wólka Mińska}
i~Feliksa\indeksosoby[Tkaczyk]{Feliksa}{Gomulska}{1897}{Desno} Gomulscy
dochowali się wspólnie siódemki dzieci:
Marianny\indeksosoby[Gomulska]{Marianna}{Świtkowska}{1917}{Desno}[Sxyzwitkowska]
(ur.~1917~r.~-- zm.~1994~r.),
Leokadii\indeksosoby[Gomulska]{Leokadia}{Karłowska}{1920}{Desno}
(ur. 1920~r. -- zm. 1976~r.), Jana\indeksosoby{Jan}{Gomulski}{1922}{Desno} (ur.
1922~r. -- zm. 1977~r.),
Stanisława\indeksosoby{Stanisław}{Gomulski}{1926}{Desno}\footnote{Stanisław
Gomulski był dziadkiem autora niniejszej książki oraz najdłużej żyjącą osobą
o~nazwisku Gomulski w~dotychczasowej historii rodziny -- przeżył dokładnie
96~lat i~333~dni.} (ur.~1926~r. -- zm. 2023~r.),
Krystyny\indeksosoby[Gomulska][Pietrzak]{Krystyna}{Banaszek}{1928}{Desno}
(ur. 1928~r. -- zm. 2012~r.),
Edwarda\indeksosoby{Edward}{Gomulski}{1935}{Desno} (ur.~1935~r.~--
zm.~2015~r.) oraz
Marcjanna\indeksosoby[Gomulska]{Marcjanna}{Główka}{1938}{Desno}\footnote{Marcjanna
Główka z~domu Gomulska żyje do~dnia dzisiejszego na~Deśnie, na~ziemi, którą
jej ojciec Józef w~1926~roku kupił wspólnie ze~swoją siostrą Karoliną,
od~rodziny Piotrkowiczów.} (ur. 1938~r.). Wszystkie dzieci Józefa i~Feliksy
Gomulskich urodziły się na~Deśnie, większość z~nich całe swoje życie związała
z~tą wsią, wyjątek stanowił jednak najmłodszy syn małżonków Gomulskich,
Edward\indeksosoby{Edward}{Gomulski}{1935}{Desno}, który po~ukończeniu szkółki
ogrodniczej przeprowadził się do~Surochowa pod Jarosławiem na~Podkarpaciu,
gdzie założył własną rodzinę.

\begin{figure}[!ht]
    \vspace*{0.5cm}
    \centering \includegraphics[width=1.0\linewidth]{
        Franciszek_Gomulski_wezwanie.jpg}
    \captionsetup{format=hang}
    \caption{Wezwanie Franciszka Gomulskiego na świadka w~rozprawie
    o~alimenty dla Marjana Zawadzkiego.}
    \label{fig:fgomulski_wezwanie}
\end{figure}

Niewiele wiadomo o~historii czwartego syna Stanisława i~Julianny Gomulskich --
\textbf{Franciszka}\indeksosoby{Franciszek}{Gomulski}{1897}{Desno}. Na~pewno
dożył wieku dorosłego, gdyż we~wrześniu 1917~roku zgłosił on~w~parafii
św.~Anny w~Długiej Kościelnej urodzenie się, wspominanego już w~bieżącym
rozdziale, Marjana Zawadzkiego\indeksosoby{Marjan}{Zawadzki}{1917}{Desno}
i~został jego ojcem chrzestnym (patrz: ryc. \ref{fig:marzawadzki_1917}).
Dożył też na~pewno do~września 1918~roku, gdyż był jednym ze świadków
w~drugiej rozprawie sądowej o~alimenty dla swojego syna chrzestnego, która
odbyła się przed Sądem Pokoju w~Otwocku 17~września 1918~roku (patrz: ryc.
\ref{fig:fgomulski_wezwanie}). Dalsze losy
Franciszka\indeksosoby{Franciszek}{Gomulski}{1897}{Desno} Gomulskiego
pozostają jednak nieznane -- nie wiadomo kiedy zmarł, czy wział ślub, albo
czy miał potomstwo.

\begin{figure}[!ht]
    \vspace*{0.5cm}
    \centering \includegraphics[width=0.75\linewidth]{
        1925_Tomasz_Gomulski_Rozalia_Gomulska_Piekut_akt_ślubu_parafia_Długa_Kościelna_wpis_2.jpg}
    \captionsetup{format=hang}
    \caption{Akt ślubu Tomasza Gomulskiego oraz Rozalii Gomulskiej z~domu
    Piekut -- par. Długa Kościelna 1925~rok (2/1925) \cite{par_dluga}.}
    \label{fig:tgomulski_1925}
\end{figure}

Piąty syn Stanisława i~Julianny Gomulskich,
\textbf{Tomasz}\indeksosoby{Tomasz}{Gomulski}{1900}{Desno}, 17~stycznia
1925~roku w~Długiej Kościelnej, wziął ślub z Rozalią
Gomulską\indeksosoby[Piekut]{Rozalia}{Gomulska}{1896}{Poręby} z~domu Piekut,
wdową po wspominanym już w~niniejszym rozdziale Janie
Gomulskim\indeksosoby{Jan}{Gomulski}{1888}{Mikanów}, synu Michała i~Marianny
Gomulskich oraz dalekim kuzynie Tomasza. Małżonkowie Gomulscy razem
wychowywali trójkę córek Rozalii z~pierwszego małżeństwa:
Franciszkę\indeksosoby[Gomulska]{Franciszka}{Gomulska}{1916}{Desno}
(ur. 1916~r. -- zm. 1984~r.), Leokadię
Mariannę\indeksosoby[Gomulska]{Leokodia Marianna}{Dobosz vel Jerzak}{1918}{Desno}
(ur.~1918~r. -- zm. ?)
i~Czesławę\indeksosoby[Gomulska]{Czesława}{Gomulska}{1922}{Desno}
(ur. 1922~r. -- zm. 2011~r.), a~także dochowali się wspólnie dwójki dzieci:
Anny\indeksosoby{Anna}{Gomulska}{1925}{Desno} (ur. 1925~r. -- zm. 2016~r.)
oraz Szczepana\indeksosoby{Szczepan}{Gomulski}{1931}{Desno} (ur. 1931~r.
-- zm. 1990~r.). Po~dziwięciu latach ich małżeństwa, w~1934~roku,
Rozalia\indeksosoby[Piekut]{Rozalia}{Gomulska}{1896}{Poręby}
zmarła, a~Tomasz Gomulski\indeksosoby{Tomasz}{Gomulski}{1900}{Desno} przez
kolejne siedem lat pozostawał wdowcem, aż~w~1941~roku, ożenił się ponownie --
ze swoją o~16~lat młodszą pasierbicą
Franciszką\indeksosoby[Gomulska]{Franciszka}{Gomulska}{1916}{Desno}, którą
wychowywał od~dziewiątego roku jej życia, najstarszą córką swojej zmarłej żony.

\begin{figure}[!ht]
    \vspace*{0.5cm}
    \centering \includegraphics[width=0.8\linewidth]{
        1941_Tomasz_Gomulski_Franciszka_Gomulska_akt_ślubu_parafia_Długa_Kościelna_wpis_10.jpg}
    \captionsetup{format=hang}
    \caption{Akt ślubu Tomasza Gomulskiego oraz Franciszki Gomulskiej -- par.
    Długa Kościelna 1941~rok (10/1941) \cite{par_dluga}.}
    \label{fig:tgomulski_1941}
\end{figure}

Co więcej Tomasz\indeksosoby{Tomasz}{Gomulski}{1900}{Desno}
i~Franciszka\indeksosoby[Gomulska]{Franciszka}{Gomulska}{1916}{Desno} byli
ze~sobą spokrewnieni -- Piotr Franciszek Gumulski, którego losy były omawiane
w~poprzednim rozdziale niniejszej książki, był pradziadkiem zarówno Tomasza,
jak i~Franciszki. To~pokrewieństwo nie było zapewne tajemnicą dla nowych
małżonków (choć mogli się o~nim dowiedzieć dopiero, gdy zaczęli planować ślub),
gdyż w~akcie ich ślubu z~dnia 9~lipca 1941~roku, zostało jasno napisane,
iż~uzyskali oni zgodę na~swój ślub od~Świętej Kongregacja ds.~Dyscypliny
Sakramentów, czyli z~samej Kurii Rzymskiej\footnote{Ze~względu na tak bliskie
powinowactwo oraz niezbyt odległe pokrewienstwo, zgody na~ich ślub nie mogła
udzielić lokalna kuria tylko musiała się na~niego zgodzić Kuria Rzymska.}
(patrz ryc. \ref{fig:tgomulski_1941}):

\begin{centeredquote}
\textit{Od powinowactwa pierwszego stopnia w~linii prostej oraz pokrewieństwa
trzeciego stopnia linii bocznej Święta Kongregacja Sakramentów udzieliła
dyspensy, reskryptem z~dnia jedenastego marca roku bieżącego za
$N^{\underline{o}}$ 3513/41.}
\end{centeredquote}

Przynajmniej jedno ze wspólnych dzieci
Tomasza\indeksosoby{Tomasz}{Gomulski}{1900}{Desno}
i~Franciszki\indeksosoby[Gomulska]{Franciszka}{Gomulska}{1916}{Desno}
Gomulskich urodziło się przed ich ślubem -- w~1937~roku na~świat przyszedł ich
najstarszy syn Karol\indeksosoby{Karol}{Gomulski}{1937}{Desno} (ur. 1937~r.~--
zm.~2003~r.). Poza Karolem, małżonkowie Gomulscy dochowali się jeszcze
co~najmmniej trójki dzieci:
Wiesławy\indeksosoby{Wiesława}{Gomulska}{1944}{Desno} (ur. 1944~r. -- zm.
2023~r.), Wawrzyńca\indeksosoby{Wawrzyniec}{Gomulski}{1945}{Desno}
(ur. 1945~r. -- zm. 2005~r.) oraz Teresy\indeksosoby{Teresa}{Gomulska}{?}{Desno}
(ur. ? -- zm. ?).

\begin{figure}[!ht]
    \vspace*{0.5cm}
    \centering \includegraphics[width=0.85\linewidth]{
        1926_Wincenty_Gomulski_Stanisława_Lesińska_akt_ślubu_parafia_Długa_Kościelna_wpis_23.jpg}
    \captionsetup{format=hang}
    \caption{Akt ślubu Wincentego Gomulskiego oraz Stanisławy Lesińskiej --
    par. Długa Kościelna 1926~rok (23/1926) \cite{par_dluga}.}
    \label{fig:wgomulski_1926}
\end{figure}

Szósty syn Stanisława i~Julianny Gomulskich,
\textbf{Wincenty}\indeksosoby{Wincenty}{Gomulski}{1903}{Desno}, 14~listopada
1926~roku w~Długiej Kościelnej, wziął ślub z~niespełna szesnastoletnią
Stanisławą Lesińską\indeksosoby[Lesińska]{Stanisława}{Gomulska}{1911}{Desno}
(ur. 1911~r. -- zm. 1986~r.), panną z~Desna, wnuczką
jednego ze współzałożycieli spółki \emph{Desna} -- Stanisława
Łabędzia\indeksosoby{Stanisław}{Łabędź}{1840}{Michałów}[Lxyzabedz]. Co~ciekwe
z~aktu ich ślubu wynika, iż~podpisali oni umowę przedślubną, 23~października
1926~roku przed notariuszem Tadeuszem Kraszewiczem w~Warszawie\footnote{Zgodnie
z~najlepszą wiedzą autora niniejszej książki, jest to~pierwszy akt ślubu
w~historii rodziny Gomulskich z~Desna, w~którym występuje informacja
o~podpisaniu umowy przedślubnej.} -- patrz ryc. \ref{fig:wgomulski_1926}.
Para dochowała się wspólnie czwórki potomstwa:
Lucyny\indeksosoby{Lucyna}{Gomulska}{1930}{Desno} (ur.~1930~r.~-- zm.~1930~r.), 
Janiny Marcjannny\indeksosoby{Janina Marcjanna}{Gomulska}{1931}{Desno}
(ur. 1931~r. -- zm. 1931~r.),
Zygmunta\indeksosoby{Zygmunt}{Gomulski}{1934}{Desno} (ur. 1934~r. -- zm.~2003~r.)
oraz Celiny\indeksosoby{Celina}{Gomulska}{1937}{Desno} (ur.~1937~r. --
zm.~1979~r.).

\begin{figure}[!ht]
    \vspace*{0.5cm}
    \centering \includegraphics[width=0.8\linewidth]{
        1940_Aleksander_Gomulski_Czesława_Gomulska_akt_ślubu_parafia_Długa_Kościelna_wpis_21.jpg}
    \captionsetup{format=hang}
    \caption{Akt ślubu Aleksandra Gomulskiego oraz Czesławy Gomulskiej --
    par. Długa Kościelna 1940~rok (21/1940) \cite{par_dluga}.}
    \label{fig:agomulski_1940}
\end{figure}

Najmłodszy syn Stanisława i~Julianny Gomulskich,
\textbf{Aleksander}\indeksosoby{Aleksander}{Gomulski}{1906}{Desno}, dnia
19~listopada 1940~roku w~Długiej Kościelnej, w~wieku 34~lat, wziął ślub
z~młodszą od niego o~16~lat Czesławą
Gomulską\indeksosoby[Gomulska]{Czesława}{Gomulska}{1922}{Desno} (ur. 1922~r. --
zm. 2011~r.), najmłodszą córką zmarłych
Jana\indeksosoby{Jan}{Gomulski}{1888}{Mikanów}
i~Rozalii\indeksosoby[Piekut]{Rozalia}{Gomulska}{1896}{Poręby} Gomulskich oraz
pasierbicą starszego brata Aleksandra --
Tomasza\indeksosoby{Tomasz}{Gomulski}{1900}{Desno}. W~związku z~tym,
że~Czesława miała w~momencie ślubu zaledwie 18~lat to~na~jej ślub musiał
wyrazić zgodę jej ojczym. Co więcej, Aleksander i~Czesława byli ze~sobą
spokrewnieni, w~takim samym stopniu jak Tomasz
i~Franciszka\indeksosoby[Gomulska]{Franciszka}{Gomulska}{1916}{Desno} (których
ślub odbył się rok później), nie było jednak między nimi bliskiego powinowactwa,
wiec wystarczyło, że~zgodę na~ich ślub wyraziła Warszawska Kuria
Metropolitarna -- dyspensę w~tej sprawie uzyskali 21~października 1940~roku
(patrz: ryc. \ref{fig:agomulski_1940}). Para dochowała się wspólnie trójki
potomstwa: Zbigniewa\indeksosoby{Zbigiew}{Gomulski}{?}{Desno} (ur. ? -- zm. ?),
Lucyny\indeksosoby{Lucyna}{Gomulski}{?}{Desno} (ur.~?~-- zm.~?) oraz
Ryszarda\indeksosoby{Ryszard}{Gomulski}{1950}{Desno} (ur. 1950~r.). Ich
najmłodszy syn Ryszard żyje do~dnia dziejszego na~Deśnie, na~ziemi, którą
jego dziadek Stanisław Gomulski odkupił w~1900~roku od Konstantnego
Kaima\indeksosoby{Konstanty}{Kaim}{1860}{Stanisławów}\footnote{Dom
Ryszarda stoi \emph{de facto} obok miejsca, w~którym stał kiedyś pierwszy dom
rodzinny Stanisława i~Julianny Gomulskich.}.

\begin{figure}[!ht]
    \vspace*{0.4cm}
    \centering \includegraphics[width=0.78\linewidth]{
        1974_Józef_Gomulski_Feliksa_Gomulska_zdjęcie_nagrobka1_cmentarz_Długa_Kościelna.jpg}
    \captionsetup{format=hang}
    \caption{Zdjęcie grobu małżonków Józefa oraz Feliksy Gomulskich,
    znajdującego się na cmentarzu parafialnym parafii św. Anny w~Długiej
    Kościelnej (sektor: A2, rząd: 8, numer: 12)}
    \label{fig:jgomulski_1974}
\end{figure}

Jak już wcześniej wspomniano, Stanisław
Gomulski\indeksosoby{Stanisław}{Gomulski}{1861}{Wólka Mińska} zmarł 3~kwietnia
1910~roku -- jego żona
Julianna\indeksosoby[Gągol]{Julianna}{Gomulska}{1862}{Dłużka} dożyła
przynajmniej do 1926~roku, gdyż była obecna na~ślubie swojego syna Wincentego.
Nie jest jednak jasne, czy dożyła do~ślubu swojego najmłodszego syna
Aleksandra\indeksosoby{Aleksander}{Gomulski}{1906}{Desno} (1940~rok) -- akt
jego ślubu tego nie precyzuje, a~aktu jej zgonu nie udało się dotąd odnaleźć.
Julianna została pochowana razem ze swoim mężem oraz najstarszym synem
w~pierwszym grobie rodzinnym rodziny Gomulskich (patrz: ryc.
\ref{fig:grob_gomulskich_1910}). Większość potomków rodziny Stanisława
i~Julianny Gomulskich spoczywa prawdopodobnie na~terenie cmentarza w~Długiej
Kościelnej, gdyż ich dzieci nie wyprowadziły się zbyt daleko
od~rodzinnej miejscowości. Ryciny: \ref{fig:jgomulski_1974},
\ref{fig:tgomulski_2008}, \ref{fig:wgomulski_1986}, \ref{fig:agomulski_2011},
\ref{fig:zgomulski_2003}, \ref{fig:szgomulski_1959}
oraz \ref{fig:sgomulski_2014} przedstawiają przykłady niektórych z~ich
nagrobków.

\begin{figure}[!ht]
    \vspace*{0.4cm}
    \centering \includegraphics[width=0.78\linewidth]{
        2008_Tomasz_Gomulski_Franciszka_Gomulska_zdjęcie_nagrobka1_cmentarz_Długa_Kościelna.jpg}
    \captionsetup{format=hang}
    \caption{Zdjęcie grobu rodzinnego Tomasza Gomulskiego, znajdującego się na
    cmentarzu parafialnym parafii św. Anny w~Długiej Kościelnej (sektor: A5,
    rząd: 2, numer: 12)}
    \label{fig:tgomulski_2008}
\end{figure}

\begin{figure}[!ht]
    \vspace*{0.4cm}
    \centering \includegraphics[width=0.78\linewidth]{
        1986_Wincenty_Gomulski_Stanisława_Gomulska_zdjęcie_nagrobka1_cmentarz_Długa_Kościelna.jpg}
    \captionsetup{format=hang}
    \caption{Zdjęcie grobu małżonków Wincentego oraz Stanisławy Gomulskich,
    znajdującego się na cmentarzu parafialnym parafii św. Anny w~Długiej
    Kościelnej (sektor: B3, rząd: 9, numer: 11)}
    \label{fig:wgomulski_1986}
\end{figure}

\begin{figure}[!ht]
    \vspace*{0.4cm}
    \centering \includegraphics[width=0.78\linewidth]{
       2011_Aleksander_Gomulski_Czesława_Gomulska_zdjęcie_nagrobka1_cmentarz_Długa_Kościelna.jpg}
    \captionsetup{format=hang}
    \caption{Zdjęcie grobu małżonków Aleksandra oraz Czesławy Gomulskich,
    znajdującego się na cmentarzu parafialnym parafii św. Anny w~Długiej
    Kościelnej (sektor: A2, rząd: 8, numer: 13)}
    \label{fig:agomulski_2011}
\end{figure}

\begin{figure}[!ht]
    \vspace*{0.4cm}
    \centering \includegraphics[width=0.78\linewidth]{
       2003_Zygmunt_Gomulski_Janina_Gomulska_zdjęcie_nagrobka1_cmentarz_Długa_Kościelna.jpg}
    \captionsetup{format=hang}
    \caption{Zdjęcie grobu małżonków Zygmunta oraz Janiny Gomulskich,
    znajdującego się na cmentarzu parafialnym parafii św. Anny w~Długiej
    Kościelnej (sektor: A4, rząd: 2, numer: 11)}
    \label{fig:zgomulski_2003}
\end{figure}

\begin{figure}[!ht]
    \vspace*{0.4cm}
    \centering \includegraphics[width=0.77\linewidth]{
       1959_Szczepan_Gomulski_zdjęcie_nagrobka_cmentarz_Długa_Kościelna.jpg}
    \captionsetup{format=hang}
    \caption{Zdjęcie grobu małżonków Szczepana oraz Władysławy Barbary
    Gomulskich, znajdującego się na~cmentarzu parafialnym parafii św.~Anny
    w~Długiej Kościelnej (sektor: A5, rząd: 2, numer: 13)}
    \label{fig:szgomulski_1959}
\end{figure}

\begin{figure}[!ht]
    \vspace*{0.4cm}
    \centering \includegraphics[width=0.77\linewidth]{
       2014_Stanisław_Gomulski_Jadwiga_Poziemska_zdjęcie_nagrobka1_cmentarz_Długa_Kościelna.jpg}
    \captionsetup{format=hang}
    \caption{Zdjęcie grobu małżonków Stanisława oraz Jadwigi Gomulskich,
    znajdującego się na cmentarzu parafialnym parafii św. Anny w~Długiej
    Kościelnej (sektor: A6, rząd: 7, numer: 19)}
    \label{fig:sgomulski_2014}
\end{figure}

\begin{figure}[!ht]
\vspace*{0.5cm}
\centering
\begin{tikzpicture}[
  % Definiujemy stałe odstępy
  node distance=1cm and 1cm, 
  every node/.style={draw, rectangle, align=center, minimum width=3.8cm,
  inner sep=5pt}]

% --- POKOLENIE 0 ---
\node (Osiowa0) [draw=none, minimum width=0pt, inner sep=0pt] at (0,0) {};
\node (Antoni) [left=1.5cm of Osiowa0] {Antoni Gomoła \\ {
    \footnotesize \emph{Wnętrzne, 1764 r.}}};
\node (Petronella) [right=1.5cm of Osiowa0] {Petronella Sokół \\ {
    \footnotesize \emph{Mysłów, 1768 r.}}};

\draw[-] (Antoni) -- (Petronella)
    node[midway, below, draw=none, minimum width=0pt,
    font=\scriptsize] (Slub0) {\emph{Wilczyska, 1783 r.}};

% --- POKOLENIE 1 ---
\node (Osiowa1) [draw=none, minimum width=0pt, inner sep=0pt,
below=1.5cm of Slub0] {};
\node (Jan) [left=1.5cm of Osiowa1] {Jan Gomoła \\ {
    \footnotesize \emph{Mysłów, 1789 r.}}};
\node (Łucja) [right=1.5cm of Osiowa1] {Łucja Rokicka \\ {
    \footnotesize \emph{Rudzienko, 1790 r.}}};

\draw[-] (Jan) -- (Łucja)
    node[midway, above, draw=none, minimum width=0pt,
    font=\scriptsize] (Slub1_Gora) { }
    node[midway, below, draw=none, minimum width=0pt,
    font=\scriptsize] (Slub1) {\emph{Kołbiel, 1814 r.}};

\draw[->] (Slub0) -- (Slub1_Gora);

% --- POKOLENIE 2 ---
\node (Osiowa2) [draw=none, minimum width=0pt, inner sep=0pt,
below=1.5cm of Slub1] {};
\node (Piotr) [left=1.5cm of Osiowa2] {Piotr Franciszek Gumulski \\ {
    \footnotesize \emph{Rudzienko, 1815 r.}}};
\node (Marianna) [right=1.5cm of Osiowa2] {Marianna Chłopik \\ {
    \footnotesize \emph{Zamienie, 1819 r.}}};
    
\draw[-] (Piotr) -- (Marianna)
    node[midway, above, draw=none, minimum width=0pt,
    font=\scriptsize] (Slub2_Gora) { }
    node[midway, below, draw=none, minimum width=0pt,
    font=\scriptsize] (Slub2) {\emph{Mińsk, 1835 r.}};

\draw[->] (Slub1) -- (Slub2_Gora);

% --- POKOLENIE 3 ---
\node (Osiowa3) [draw=none, minimum width=0pt, inner sep=0pt,
below=1.5cm of  Slub2] {};
\node (Jan2) [left=1.5cm of Osiowa3] {Jan Gomulski \\ {
    \footnotesize \emph{Zamienie, 1836 r.}}};
\node (Agnieszka) [right=1.5cm of Osiowa3] {Agnieszka Piotrkowicz \\ {
    \footnotesize \emph{Budy Barcząckie, 1840 r.}}};

\draw[-] (Jan2) -- (Agnieszka)
    node[midway, above, draw=none, minimum width=0pt,
    font=\scriptsize] (Slub3_Gora) { }
    node[midway, below, draw=none, minimum width=0pt,
    font=\scriptsize] (Slub3) {\emph{Mińsk, 1858 r.}};

\draw[->] (Slub2) -- (Slub3_Gora);

% --- POKOLENIE 4 ---
\node (Osiowa4) [draw=none, minimum width=0pt, inner sep=0pt,
below=1.5cm of  Slub3] {};
\node (Stanislaw) [left=1.5cm of Osiowa4] {Stanisław Gomulski \\ {
    \footnotesize \emph{Wólka Mińska, 1861 r.}}};
\node (Julianna) [right=1.5cm of Osiowa4] {Julianna Gągol \\ {
    \footnotesize \emph{Dłużka, 1862 r.}}};

\draw[-] (Stanislaw) -- (Julianna)
    node[midway, above, draw=none, minimum width=0pt,
    font=\scriptsize] (Slub4_Gora) { }
    node[midway, below, draw=none, minimum width=0pt,
    font=\scriptsize] (Slub4) {\emph{Mińsk, 1881 r.}};

\draw[->] (Slub3) -- (Slub4_Gora);

\end{tikzpicture}
\end{figure}

Bezpośredni potomkowie
Stanisława\indeksosoby{Stanisław}{Gomulski}{1861}{Wólka Mińska}
i~Julianny\indeksosoby[Gągol]{Julianna}{Gomulska}{1862}{Dłużka} Gomulskich,
mieszkają na~terenie obecnego Desna do~dnia dzisiejszego -- część z~nich nosi
nazwisko Gomulski, ale poza nimi, Desno zamieszkują również rodziny
spokrewnione o~nazwiskach: Główka, Kochaniak, Banaszek, Dybek, Gryz,
Rupniewski, Szatański, Marczak czy Karłowski, które również wywodzą się
z~tej, najmłodszej na~Deśnie, gałęzi rodziny Gomulskich.